\section{Tropical \texorpdfstring{$2$}{2}-cycles on the small resolution}\label{sec:resolution}

This section proves Theorem~\ref{thm:res}: for every tropical $2$-cycle $(S,j,v)$ on $B'_\sigma$ in the sense of
\cite[Def.~7.2]{CBM} we construct an integral $3$-chain $\Xi(S)$ on a projective small resolution $\widehat X$ of
$X'_\sigma$ with $\partial\Xi(S)=\sum_q\rho_q(S)[C_q]$, and we compare the other small resolutions. It is the
resolution-side counterpart of Theorem~\ref{thm:lift} and forms Part~III of the paper. No result of this section is used
in Parts~I and~II: Theorems~\ref{thm:main}, \ref{thm:real} and~\ref{thm:lift} and Corollary~\ref{cor:smooth} are proved
without it; the integral relations among the exceptional curves are established independently in
Proposition~\ref{prop:integral}. Conversely, this section uses the following results of the other parts. From Section~\ref{sec:base}:
Proposition~\ref{prop:partialres} and Lemma~\ref{lem:H2} on the small resolutions, Remark~\ref{rem:tie} on the ties, and,
on the base, Lemmas~\ref{lem:monodromy}, \ref{lem:stack} and~\ref{lem:types} and the passages at a node. From
Section~\ref{sec:cycles}: the clauses (a)--(j) of \cite[Def.~7.2]{CBM} and Lemma~\ref{lem:boundaryfield}
(Section~\ref{ssec:cbm}), Lemma~\ref{lem:lift}, the node coefficients (Section~\ref{ssec:cycles-bdry}) and, in a remark,
Proposition~\ref{prop:containment}. From Section~\ref{sec:dualpair}: Propositions~\ref{prop:L1} and~\ref{prop:L2} on the
dual-pair complex, whose proofs do not use the present section. Theorems~\ref{thm:real} and~\ref{thm:lift}, and
Proposition~\ref{lem:F}, enter only the last assertion of Theorem~\ref{thm:res} and Section~\ref{ssec:res-relative}.
Of the hypotheses we use~\ref{hyp:P} and the pair $(\mathcal T,\check h)$ of Theorem~\ref{thm:U}, with $\check h$
normalised as in Section~\ref{ssec:gross}, so that $\check h(m)$ is odd and positive for every
$m\in\partial\Delta^\circ\cap M$; the pulling refinement and the mirror family play no role. Throughout, $j$ is piecewise
linear.

\subsection{Outline of the construction}\label{ssec:res-outline}
Casta\~no-Bernard and Matessi attach to a tropical $2$-cycle $S$, in their topological models, a three-dimensional object
$\widetilde S^*$ \cite[\S7]{CBM}: over the smooth part of $S$ it is the bundle of circles $\R v/\Z v$ in the fibres of the
torus fibration, closed off over $\partial S$, where the circle collapses, and completed over the interior edges and vertices
of $S$ by triangles and tetrahedra \cite[\S7.5, \S7.7]{CBM}. We realise this object in the
algebraic threefold $\widehat X$. The chain is $\Xi(S)=(\Upsilon_{\rm res})_*\Xi_W$, where $\Xi_W$ is, up to sign and up to chains
inside the exceptional curves, the image $F_*[\mathfrak D(S)]$ of the fundamental chain of a cell complex $\mathfrak D(S)$
under a continuous map $F$ into a model $W$ of $\widehat X$, and $\Upsilon_{\rm res}$ is a homeomorphism of the ambient toric
variety carrying $W$ onto $\widehat X$. The construction has five steps.

\emph{The resolved base} (Sections~\ref{ssec:res-base} and~\ref{ssec:res-shell}). The resolution $\widehat X$ is induced by a
regular refinement $\mathcal F'$ of $\Sigma'$ that splits every node parallelogram along a diagonal. Gross's construction
with $\mathcal F'$ in place of $\Sigma'$ gives a second base $\widehat B$ on the same sphere $\partial\nabla^{\check h}$, in
which each vertex $p_q$ of the discriminant is replaced by two negative vertices $y_0,y_\infty$ joined by an arc $e_q$ of the
discriminant. After $S$ is put in a standard position near the nodes, a homeomorphism $B'_\sigma\to\widehat B$ carries it
to a surface $\hat S$, the \emph{resolved cycle}, that satisfies the clauses of \cite[Def.~7.2]{CBM} away from the nodes, while near each node met
by $\partial S$ its boundary turns at $y_0$ and $y_\infty$ and runs along $e_q$, whose vanishing vector is not the field
of $S$. This configuration produces the exceptional curve $C_q$.

\emph{The localized hypersurface} (Section~\ref{ssec:res-W}). The model $W$ is the localization, in the sense of Mikhalkin
\cite{Mikhalkin} and Yamamoto \cite{Ya16}, of a member $X_R$ of Gross's degeneration: each monomial is cut off where it
is far from the maximal one. Three properties of $W$ are used. Near every torus orbit that it meets, $W$ is given by an
exact equation in the active monomials, and it is isotopic to the proper transform of $X_R$ through hypersurfaces
transverse to the orbits. Its positive real points project homeomorphically onto $\widehat B$, which gives a section
$\mathfrak s\colon\widehat B\to W$. Over small open sets $U\subset\widehat B$ adapted to the faces of $\nabla^{\check h}$ and
disjoint from the discriminant, the preimages $W_U$ are aspherical, with fundamental group the lattice of integral tangent
vectors of $\widehat B$, and inclusions induce parallel transport. Consequently the only obstruction to extending a map
over the cells of $\mathfrak D(S)$ above the regular part of $S$ is the balancing condition of $S$, which holds.

\emph{The domain and the local pieces} (Sections~\ref{ssec:res-domain} and~\ref{ssec:res-local}). The complex
$\mathfrak D(S)$ is the circle bundle above, completed near $\partial S$ and near the points where $S$ crosses the
discriminant by explicit local domains. The positive real points of $W$ avoid every toric stratum, whereas the circle of
$v$ collapses only on a stratum; so near $\partial S$ the pieces are swept out by the circle of $v$ acting on
\emph{collar paths}, explicit paths in $W$ from the positive section to the critical locus of the local toric model,
along which the phase of one character turns by $\pi$. The collar paths of neighbouring points of $\partial S$ must be
joined by discs in $W$, and the obstruction is the total change of the phases of finitely many characters around the
loop they form (Lemma~\ref{lem:pi1}). At the nodes and at the negative vertices it vanishes because all collar paths
there lie in one two-face and turn the same character by the same angle. At the positive vertices it vanishes after
\emph{slides}, which rotate the collar paths of one leg by an integral multiple of $\pi$ in a circle that preserves $W$;
where a sheet changes side of a two-face, the two families of collar paths are joined across a \emph{switching disc}; and at
a node on the negative side the filling is made explicit by a half-turn along the exceptional curve (Lemma~\ref{lem:collar}).
On the mirror side the corresponding pieces fit together without these corrections (Section~\ref{sec:lift}), because there
the fibre over a sheet is the two-torus $\ker v\otimes U(1)$ rather than a circle.

\emph{The boundary.} The boundary of $\mathfrak D(S)$ consists of one $2$-sphere $\partial_q\mathfrak D$ for each node on
$\partial S$, the circle bundle over $e_q$ collapsed at its ends; the node piece maps it onto the exceptional curve of $W$
with degree $-\rho_q(S)$ (Lemmas~\ref{lem:domain} and~\ref{lem:nodedeg}).

\emph{Transport} (Section~\ref{ssec:res-transport}). A flow tangent to the torus orbits carries $W$ to $\widehat X$ and its
exceptional curves to the curves $C_q$ with degree one (Lemmas~\ref{lem:T0} and~\ref{lem:T}).

The construction involves choices, made in the order listed at the beginning of Section~\ref{ssec:res-proof}; the chain
$\Xi(S)$ is determined by them up to closed chains, and its boundary does not depend on them. The combinatorial data at
the two-faces and legs of $\nabla^{\check h}$ come first. Next, $S$ is put in standard position near the nodes and
moved to the resolved base $\widehat B$, and its germ along $\partial S$ is normalised. Then a neighbourhood $\mathcal N_\Gamma$ of
$\partial S$ and of the crossing points, a triangulation of the rest of $S$ and the associated open sets of $\widehat B$ are
fixed. The cut-off constants and the parameter $R$ of $W$ follow, then the local pieces and the extension over the rest of
$S$, and finally the transport to $\widehat X$. Five technical proofs are deferred to Appendix~\ref{app:original}.

\subsection{The resolved base}\label{ssec:res-base}
We compare the discriminants before and after splitting each node parallelogram along a diagonal.
The dual-pair description controls both the topology of this modification and the transport of the integral tangent lattices.

We use the dual-pair complex of Section~\ref{sec:dualpair}. For a fan-side subdivision $\mathcal F$ of $\partial\Delta$, a
pair $(\tau,\omega)$ of a simplex $\tau$ of $\mathcal T$ and a cell $\omega$ of $\mathcal F$ is \emph{dual} if
$\langle u,n\rangle=-1$ for all vertices $u$ of $\tau$ and $n$ of $\omega$. The base built from $(\mathcal T,\check h)$ and
$\mathcal F$ is the disjoint union of open cells $C_B(\tau,\omega)$ of dimension $3-\dim\tau-\dim\omega$, one for each dual
pair, whose closures are balls (Proposition~\ref{prop:L1}); the discriminant is the union of the cells with $\dim\tau\ge1$
and $\dim\omega\ge1$, a graph whose edges, the \emph{arcs}, are the one-cells $C_B(\tau,\omega)$ with $\tau$ and $\omega$
edges, each a union of legs (Proposition~\ref{prop:L2}). We write $F_\tau$ for the face of $\nabla^{\check h}$ with inner
normal cone $\operatorname{cone}(\tau)$, and $F_u$, $F_{uu'}$ for its facets and two-faces.

\emph{Labels.} In this section the diagonal of a node parallelogram $P=abcd$ is always the diagonal $\{a,c\}$ that the
resolution $\widehat X$ under consideration induces on it; by Section~\ref{ssec:admissible} this fixes the circuit
$\circv_q$ and the node coefficients $\rho_q$, and a change of diagonal reverses both signs. Projective small resolutions
exist, with the diagonal through the centre at every node parallelogram of a pentagon or a hexagon
(Proposition~\ref{prop:partialres}(3)). What this section needs beyond Proposition~\ref{prop:partialres}(3) is that the
refinement can be taken not to subdivide anything else in the faces of dimension at most two.

\begin{lemma}\label{lem:regref}
Let $\widehat X$ be a projective small resolution of $X'_\sigma$. There is a regular refinement $\mathcal F'$ of $\Sigma'$, with
rays generated by lattice points of $\partial\Delta$, which induces the diagonal of $\widehat X$ on every node parallelogram
and subdivides no other cell of $\Sigma'$ lying in a face of $\Delta$ of dimension at most two; the proper transform of $X$ in
the toric variety of $\mathcal F'$ is $\widehat X$.
\end{lemma}
\begin{proof}
Apply Proposition~\ref{prop:partialres}(3) with the diagonals of $\widehat X$ as the chosen diagonals. Since $\widehat X$ is
projective, no nonzero $\lambda\ge0$ lies in $K$, and the proof of Proposition~\ref{prop:partialres}(3) produces a function
$\psi$ with $\langle\circv_q,\psi\rangle>0$ for every $q$ and the refinement induced by $h+\epsilon\psi$, $h$ strictly convex
inducing $\Sigma'$ and $\epsilon>0$ small. It subdivides each cell $\omega$ of $\Sigma'$ by the regular subdivision of
$\psi|_\omega$. The cells of $\Sigma'$ in faces of dimension at most two are unimodular triangles, primitive segments and node
parallelograms (Definition~\ref{def:profile}); the first two have no lattice points besides their vertices and are not
subdivided, and each node parallelogram receives the diagonal $\{a,c\}$ of $\widehat X$, since $\langle\circv_q,\psi\rangle>0$.
The last assertion is part of Proposition~\ref{prop:partialres}(3).
\end{proof}

\emph{Throughout Sections~\ref{ssec:res-base}--\ref{ssec:res-transport}, $\widehat X$ is projective}, and we fix
$\mathcal F'$ as in Lemma~\ref{lem:regref}, generic inside the facets of $\Delta$, with toric variety $X_{\mathcal F'}$. For
a small resolution that is not projective the chain is constructed on a projective one and carried over by
Lemma~\ref{lem:flop}. Let $\widehat B$ be Gross's base (Section~\ref{ssec:gross}) built from the same pair
$(\mathcal T,\check h)$ and the fan side $\mathcal F'$ in place of $\Sigma'$, with the subdivision of the upper and lower
points given by heights $h'$ on the lower points and $0$ on the upper points, $h'$ a strictly convex function inducing
$\mathcal F'$ (Section~\ref{sec:dualpair}). It is the same polyhedral sphere $\partial\nabla^{\check h}$ as $B'_\sigma$,
with the same facets and facet charts; only the fan side differs. We write $\Gamma$ for the discriminant of $B'_\sigma$ and
$\widehat\Gamma$ for that of $\widehat B$, and call $\widehat B$ the \emph{resolved base}. It is the base of Gross's toric
degeneration of $\widehat X$, just as $B'_\sigma$ is that of $X'_\sigma$.

\emph{The node in $\widehat B$.} Let $q\in Q$, with parallelogram $P=abcd$ in the two-face $G$ of $\Delta$, and let
$G^\circ=[u_-,u_+]$; the vertex $p_q$ of the discriminant lies in the two-face $F_q:=F_{u_+u_-}$ of $\nabla^{\check h}$. In $\mathcal F'$
the parallelogram is replaced by the unimodular triangles $\{a,b,c\}$ and $\{a,c,d\}$, which by Lemmas~\ref{lem:stack}
and~\ref{lem:types} give negative vertices $y_0$ and $y_\infty$ of $\widehat\Gamma$ in $F_q$. They are joined by the arc
$e_q=C_{\widehat B}([u_+,u_-],[a,c])$ (Proposition~\ref{prop:L2}), a union of legs with lower support $\{a,c\}$ and vanishing
vector $c-a$. The legs with lower supports $\{a,b\}$ and $\{b,c\}$ end at $y_0$, those with $\{c,d\}$ and $\{d,a\}$ at
$y_\infty$.

\begin{lemma}[two-faces]\label{lem:twofaces}
Let $B$ be $B'_\sigma$ or $\widehat B$, with discriminant $\Gamma_*$ and cells $C_B(\tau,\omega)$. Let $F_{uu'}$ be a
two-face of $\nabla^{\check h}$, $\tau=[u,u']$, and $G=\tau^\vee$ the dual face of $\Delta$; then $\dim G\le2$, and every
lattice point of $G$ is a vertex of the subdivision of $G$.
\begin{enumerate}
\item $\operatorname{ri}F_{uu'}\setminus\Gamma_*$ is the disjoint union of the open discs
$\widetilde R_n\cap\operatorname{ri}F_{uu'}=C_B(\tau,n)$, $n\in G\cap N$.
\item $\Gamma_*\cap\operatorname{ri}F_{uu'}$ is the union of the arcs $C_B(\tau,[n,n'])$, $[n,n']$ an edge of the subdivision
of $G$, and of the vertices $C_B(\tau,\omega)$, $\omega$ a two-cell of the subdivision. The arc $C_B(\tau,[n,n'])$ lies in the
closures of $C_B(\tau,n)$ and $C_B(\tau,n')$ and in no other closure $\overline{C_B(\tau,n'')}$.
\item The closure of $C_B(\tau,n)$ is a closed piecewise linear disc, the union of the cells $C_B(\tau',\omega)$ with
$\tau'\supseteq\tau$ and $n\in\omega$.
\end{enumerate}
Thus $\Gamma_*\cap F_{uu'}$ with its regions is the dual-block decomposition of $G$ with its subdivision, together with
the arcs running to $\partial F_{uu'}$. For an edge $F_{\tau'}$ of $\nabla^{\check h}$ the same holds with the open intervals
$C_B(\tau',n)$ and the points $C_B(\tau',[n,n'])$.
\end{lemma}
Here $\widetilde R_n$ is the vertex region of $n$ (Section~\ref{ssec:gross}), $\operatorname{ri}$ denotes the relative interior,
and the lattice points of $G$ are its vertices and, for a pentagon or a hexagon, its centre.
\begin{proof}
Apply Propositions~\ref{prop:L1} and~\ref{prop:L2} with $\mathcal F=\Sigma'$ and with $\mathcal F=\mathcal F'$. By
Proposition~\ref{prop:L1}(5), $\operatorname{ri}F_{uu'}$ is the union of the cells $C_B(\tau,\omega)$, and
$\widetilde R_n\cap\operatorname{ri}F_{uu'}=C_B(\tau,n)$, which is nonempty if and only if $(\tau,n)$ is a dual pair, that
is, $n\in G\cap N$. By Proposition~\ref{prop:L2}(1), $\Gamma_*$ is the union of the cells $C_B(\tau',\omega)$ with
$\dim\tau'\ge1$ and $\dim\omega\ge1$; intersecting with $\operatorname{ri}F_{uu'}$ gives the cells $C_B(\tau,\omega)$ with
$\dim\omega\ge1$, which are the arcs ($\omega$ an edge) and the vertices ($\omega$ a two-cell) of (2), and their complement
in $\operatorname{ri}F_{uu'}$ is the union of the $C_B(\tau,n)$. By Proposition~\ref{prop:L1}(2), $C_B(\tau,n)$ is an open
cell of dimension $3-\dim\tau-0=2$, and its closure is the piecewise linear disc $\bigcup_{p\preceq(\tau,n)}C_B(p)$; since
$(\tau',\omega)\preceq(\tau,n)$ means $\tau'\supseteq\tau$ and $n\in\omega$, this gives (3) and the incidence statement of
(2). The case of an edge is the same with $\dim\tau'=2$.
\end{proof}

Since $\Sigma'$ and $\mathcal F'$ induce the same subdivision on every face of $\Delta$ of dimension at most two except at
the node parallelograms, Lemma~\ref{lem:twofaces} gives isomorphic cell structures, with the same vertex regions, on
every two-face of $B'_\sigma$ and $\widehat B$ away from the node points. Near a node, let $U_q$ be the open star of $p_q$ in
the barycentric subdivision of the polyhedral decomposition $\mathscr P$ of $B'_\sigma$. It lies in
$\operatorname{ri}F_q\cup\operatorname{int}F_{u_+}\cup\operatorname{int}F_{u_-}$, and it is the open cone from $p_q$ over its
link $\operatorname{lk}_q\cong S^2$, linear on each simplex; for $0<r\le1$ let $U^r_q$ be the part of cone parameter $<r$.
The set $\Gamma\cap U_q$ consists of the four legs at $p_q$, which are radial and end at points
$\ell_{ab},\ell_{bc},\ell_{cd},\ell_{da}$ of $\operatorname{lk}_q$, in this cyclic order on the circle
$\operatorname{lk}_q\cap F_q$; here $\ell_{ab}$ is the end of the leg with lower support $\{a,b\}$. The simplices of the
barycentric subdivision correspond to chains of cells of $\mathscr P$, and $U_q$ is the union of the open simplices whose
chain contains the carrying cell of $q$; since no chain contains two different two-cells, distinct $U_q$ are disjoint. For
$r<1$ the closures of the $U^r_q$ are therefore pairwise disjoint as well. This holds even when two closed carrying cells
share an edge, as happens at a hexagon on the $A_2$ branch.

\begin{lemma}[moving to the resolved base]\label{lem:resmove}
For every $r_*\in(0,1)$ there is a piecewise linear homeomorphism $g$ of $\partial\nabla^{\check h}$ with the following
properties.
\begin{enumerate}
\item $g$ preserves every closed face of $\nabla^{\check h}$ and is the identity on each facet interior outside a collar of
its boundary.
\item $\widehat U_q:=g(U_q^{r_*})$ is an open ball containing $y_0\cup e_q\cup y_\infty$, and $g$ maps
$\Gamma\setminus\bigcup_qU_q^{r_*}$ onto $\widehat\Gamma\setminus\bigcup_q\widehat U_q$, the part of each leg at $p_q$
outside $U^{r_*}_q$ onto the outer part of the leg of $\widehat B$ with the same lower support, and, on the two-skeleton
$\bigcup F_{uu'}$ outside $\bigcup_qU^{r_*}_q$, each vertex region $\widetilde R_n$ into the vertex region of the same $n$.
\item Off $\bigcup_qU_q^{r_*}$, $g$ carries the local system $\Lambda$ of $B'_\sigma$ to that of $\widehat B$, by the
identity on each facet lattice $u^\perp\cap N$, and it preserves the types of the vertices of the discriminant.
\end{enumerate}
\end{lemma}
\begin{proof}
We build $g$ skeleton by skeleton.

\emph{Edges.} On each edge of $\nabla^{\check h}$ the two discriminants meet the edge in finite sets separating intervals in
the same vertex regions, in the same order (Lemma~\ref{lem:twofaces} for the edge), and $g$ is the piecewise linear map
matching them.

\emph{Two-faces without node points.} On a two-face that contains no node point, the two cell structures, with their vertex
regions, are isomorphic, compatibly with $g$ on the boundary; $g$ extends over each closed region, a closed disc
(Lemma~\ref{lem:twofaces}(3)), by the Alexander trick \cite[Lemma~2.1]{FM12}.

\emph{Two-faces with node points.} Let $F$ be a two-face containing node points, and $G$ the dual face of $\Delta$. It
contains one node point if $G$ is a square, two if $G$ is a pentagon or a hexagon on the $A_1$ branch, and three if $G$ is a
hexagon on the $A_2$ branch. For each node point $p_q\in F$ let $F'_q=U^{r_*}_q\cap F$, an open disc around $p_q$ whose
boundary meets the four legs in the cyclic order $ab,bc,cd,da$, and let $\widehat F_q$ be a closed regular neighbourhood in
$F$ of the tree $y_0\cup e_q\cup y_\infty$ that meets $\widehat\Gamma$ only in the tree and four initial leg segments. Its
boundary meets these legs in the same cyclic order; the region of $b$ lies between $ab$ and $bc$, that of $d$ between $cd$ and
$da$, and the regions of $a$ and $c$ each border $e_q$.

The closures $\overline{F'_q}$ are pairwise disjoint, since $r_*<1$. The trees are pairwise disjoint as well: the tree of $q$
is the union of the cells $C_{\widehat B}([u,u'],\omega)$ with $\omega$ one of the two triangles into which the diagonal of
$\widehat X$ cuts the parallelogram of $q$, or that diagonal, and different parallelograms give different cells, which are
disjoint (Proposition~\ref{prop:L1}). With the diagonals through the centre $o$ this says that the diagonals are pairwise
non-adjacent edges of the star subdivision of $G$, so that no triangle of the star belongs to two trees; at a hexagon on the
$A_2$ branch the three trees use all six triangles, and consecutive trees are joined by the arcs across the edges
$[o,w_j]$, $j\equiv j_0\pmod2$. We take the neighbourhoods $\widehat F_q$ so small that they are pairwise disjoint; each of them
then meets such a joining arc only in an initial segment.

Since $\Gamma\cap U_q$ consists of the four radial legs, $F'_q$ is the open cone from $p_q$ over the circle
$\operatorname{lk}_q\cap F$, and it meets the closure of each region in the cone over an arc of that circle: a sector bounded by
two radial leg segments, those of $ab$ and $bc$ for the region of $b$, of $bc$ and $cd$ for that of $c$, and so on. Likewise
$\widehat F_q$, a regular neighbourhood of a tree lying on the boundaries of the regions, meets the closure of each region in a
disc attached along an arc of its boundary.

By Lemma~\ref{lem:twofaces}, $\Gamma\cap F$ and $\widehat\Gamma\cap F$ are the dual-block graphs of the two subdivisions of
$G$, together with the arcs to $\partial F$: a two-cell $\omega$ gives the vertex $C_B([u,u'],\omega)$, and an edge of the
subdivision gives the arc $C_B([u,u'],[n,n'])$ joining the vertices of the two-cells that contain it, or running to $\partial F$
if only one does. The subdivision induced by $\mathcal F'$ differs from that of $\Sigma'$ only by the diagonal $\{a,c\}$ in each
node parallelogram. So the graph of $\widehat B$ is obtained from that of $B'_\sigma$ by replacing each four-valent vertex
$p_q$ by the tree $y_0\cup e_q\cup y_\infty$, where $y_0$ (the triangle $\{a,b,c\}$) carries the arcs across $[a,b]$ and
$[b,c]$, and $y_\infty$ (the triangle $\{a,c,d\}$) those across $[c,d]$ and $[d,a]$; every other vertex, every other arc and
every incidence with a region is unchanged. After removing the discs $F'_q$ from $F$ and the neighbourhoods $\widehat F_q$,
the two labelled graphs, with their regions, are therefore isomorphic, the four boundary points of each removed disc being
matched in the cyclic order $ab,bc,cd,da$. Each closed region $\overline{C_B([u,u'],n)}$ is a closed disc
(Lemma~\ref{lem:twofaces}(3)), and by the previous paragraph the removed discs meet it in pairwise disjoint discs attached
along pairwise disjoint arcs of its boundary, so what remains of it is again a closed disc, on both bases. Hence $g$ extends over $F\setminus\bigcup_qF'_q$ region by region, by the
Alexander trick, and over each $F'_q$ onto $\operatorname{int}\widehat F_q$ by coning.

The shape of the graph inside $F$ depends on $G$. For a square it has one four-valent vertex. For a pentagon it has two
four-valent vertices, joined by the arc across the edge that the two parallelograms share, and one negative vertex, the
triangle, joined to both. For a hexagon on the $A_1$ branch it has two four-valent vertices and two negative vertices,
forming a cycle of four arcs around the region of the centre $o$. For a hexagon on the $A_2$ branch the three node points are
joined pairwise by three arcs across the edges $[o,w_j]$, $j\equiv j_0\pmod 2$, which bound the region of $o$ in $F$ and form a
cycle without trivalent vertices; on $\widehat B$ the six triangles of the star subdivision of $G$ give a cycle of six negative
vertices, three of whose arcs are the $e_q$. In every case the argument above applies unchanged, since it only uses the local
replacement at each node point.

\emph{Facets.} On the boundary of a facet $F_u$, a piecewise linear two-sphere, $g$ fixes every vertex of
$\nabla^{\check h}$ and maps every edge and every two-face of $F_u$ onto itself; in particular it is an orientation-preserving
piecewise linear homeomorphism of $\partial F_u$. We extend it over a collar by a piecewise linear homeomorphism $\mathcal G$ of
$\partial F_u\times[0,1]$ that equals $g$ on $\partial F_u\times\{1\}$, the identity on $\partial F_u\times\{0\}$, and maps
$\phi\times[0,1]$ onto itself for every face $\phi$ of $\partial F_u$; it is built on these prisms by increasing dimension. On
$v\times[0,1]$, $v$ a vertex, $\mathcal G$ is the identity. On $e\times[0,1]$, $e\cong[0,1]$ an edge, it is the Alexander
trick: $\mathcal G(x,t)=(t\,g(x/t),t)$ for $x\le t$ and $\mathcal G(x,t)=(x,t)$ for $x\ge t$. The set $\{x\le t\}$ is the cone
from $(0,0)$ over $e\times\{1\}$, and on the cone over a segment of $e\times\{1\}$ on which $g(x)=\alpha x+\beta$ the map is
$(x,t)\mapsto(\alpha x+\beta t,t)$, which is linear; so $\mathcal G$ is a piecewise linear homeomorphism of the square, and it
is the identity on the vertical sides. On $D\times[0,1]$, $D$ a two-face, $\mathcal G$ is now defined on the boundary sphere
$\partial(D\times[0,1])$ as a piecewise linear homeomorphism onto itself, and we extend it by coning from an interior point $c_D$: if the boundary map is simplicial from a
triangulation $K$ of the sphere onto a triangulation $L$, the cone map is the simplicial isomorphism from the cone $c_D*K$ onto
$c_D*L$, both of which triangulate the convex cell $D\times[0,1]$. Finally let $c$ be an interior point of $F_u$ and $0<\lambda<1$, triangulate the prisms
$\phi\times[0,1]$ compatibly and without new vertices, and map each vertex $(w,t)$ to $c+(\lambda+(1-\lambda)t)(w-c)$ and each
simplex linearly. As in the proof of Lemma~\ref{lem:NS} this is a piecewise linear homeomorphism of $\partial F_u\times[0,1]$
onto the collar between $\partial F_u$ and $c+\lambda(\partial F_u-c)$. Transport $\mathcal G$ to the collar by it, and put $g$
equal to the identity on $c+\lambda(F_u-c)$; this gives (1). Since $U_q$ meets the two-skeleton only in $\operatorname{ri}F_q$, $\widehat U_q$ is an open ball with
$\widehat U_q\cap F_q=\operatorname{int}\widehat F_q$, and since $\widehat\Gamma$ lies in the two-skeleton,
$\widehat\Gamma\cap\widehat U_q=\widehat\Gamma\cap\operatorname{int}\widehat F_q$.

\emph{Property (3).} Both local systems are the facet lattices $u^\perp\cap N$ on facet interiors, and the transition across
a two-face at a point of the vertex region of $n$ is $u^\perp\cap N\to N/\Z n\to u'^\perp\cap N$ on both bases; $g$ preserves
the open facets, and on the two-skeleton outside the node discs it preserves the vertex regions. Vertex types are determined
by the lower supports (Lemma~\ref{lem:types}), which $g$ preserves.
\end{proof}

\emph{Del Pezzo faces.} The construction is local at each node, and the global shape of a two-face $G$ of $\Delta$ enters
only at two places: Lemma~\ref{lem:resmove}, where a two-face of $\nabla^{\check h}$ may carry two or three node points, and
Remark~\ref{rem:res-roots}, which describes the boundary at a del Pezzo point by means of the ties of Remark~\ref{rem:tie}. Everything
else uses only that a node parallelogram is a unit parallelogram in a two-face $G$ with $\ell(G^\circ)=1$, that the other cells
of $\Sigma'$ in $G$ are unimodular triangles and primitive segments, and that distinct node points have disjoint stars.

\subsection{Standard position near the nodes}\label{ssec:res-shell}
Near a node the germ of $S$ may wind around the legs. In a spherical shell around $p_q$ this winding can be traded for
finitely many ordinary crossings of $\Gamma$, and inside the shell $S$ then has a standard form.

\begin{lemma}[standard position near the nodes]\label{lem:NS}
Let $(S,j,v)$ be a tropical $2$-cycle on $B'_\sigma$, with $j$ piecewise linear. There are $r_*\in(0,1)$ and a tropical
$2$-cycle $(S,j',v)$ such that:
\begin{enumerate}
\item[(0)] $j'=j$ outside $\bigcup_qU_q^{2r_*}$, and $j'(\partial S)=j(\partial S)$ with the same legs at every node and the
same boundary fields; hence $\rho(S,j')=\rho(S,j)$;
\item[(1)] for each node $q$, $j'(S)\cap U^{r_*}_q=\emptyset$ if $p_q\notin j(\partial S)$; otherwise $j'(S)\cap U^{r_*}_q$ is
the cone from $p_q$ over a piecewise linear arc $\gamma^0_q\subset\operatorname{lk}_q$ which joins the ends of the two
opposite legs of the passage and lies, apart from its ends, in one of the two open hemispheres of
$\operatorname{lk}_q\setminus F_q$, on a side $s_q\in\{\pm1\}$;
\item[(2)] $U^{r_*}_q$ contains no crossing point $\mathfrak x_k$, no interior edge or vertex of $S$, and no other sheet;
\item[(3)] the closures of the $U^{r_*}_q$ are pairwise disjoint.
\end{enumerate}
\end{lemma}
Here the passages at a node are those of Sections~\ref{ssec:disc} and~\ref{ssec:cbm}: by clause (e), $\partial S$ follows either the legs
with lower supports $\{a,b\}$ and $\{c,d\}$ (the \emph{$b-a$ passage}) or those with $\{b,c\}$ and $\{d,a\}$ (the
\emph{$c-b$ passage}). The side $s_q=+1$ is that of the facet $F_{u}$ and $s_q=-1$ that of $F_{u'}$,
for the ordering $(u,u')$ of $F_q$ fixed in Section~\ref{ssec:res-local}.
\begin{proof}
Choose triangulations of $S$ and of $B'_\sigma$, the latter subdividing the barycentric subdivision near each $p_q$ and having
$p_q$ as a vertex, for which $j$ is simplicial. Since the star of a vertex is the cone over its link, for small $r_0$ the set
$j(S)\cap U^{r_0}_q$ is the union of the radial segments through $\gamma=j(S)\cap\partial U^{r_0}_q$. If $p_q\notin j(S)$
this set is empty for small $r_0$. Otherwise $p_q=j(p)$ with $p\in\partial S$ a smooth boundary point, since a four-valent
point of $\Gamma$ is neither the image of a crossing point $\mathfrak x_k$ nor a negative vertex (clauses (b), (d)). The
crossing points, the interior edges and vertices, and $j(S\setminus V)$ for a neighbourhood $V$ of $p$ are compact and miss
$p_q$, so for small $r_0$ the set $j(S)\cap U_q^{r_0}$ is the germ of $j(S)$ at $p_q$, and by clause (e) $\gamma$ is an arc from
$\ell_{ab}$ to $\ell_{cd}$ (the $b-a$ passage) or from $\ell_{bc}$ to $\ell_{da}$ (the $c-b$ passage).

Put $r_*=r_0/2$ and let $\gamma^0_q$ be an arc on side $s_q$ with the same ends. Arcs in $S^2$ with fixed ends are
isotopic; choose a piecewise linear isotopy $\gamma_r$, $r\in[r_*,r_0]$, from $\gamma^0_q$ to $\gamma$ fixing the ends and
transverse to the other two leg ends, which it therefore crosses finitely often. Define $j'$ on $U^{r_0}_q$ as the cone over
$\gamma^0_q$ for $r\le r_*$ and as $\gamma_r$ at level $r$ for $r\in[r_*,r_0]$. Different levels are disjoint, so $j'$ is an
embedding, and it agrees with $j$ near $\partial U^{r_0}_q$. The map $(x,r)\mapsto rx$ from
$\operatorname{lk}_q\times[r_*,r_0]$ to the shell is not piecewise linear, so we realise the levels differently: triangulate
the prisms $\sigma\times[r_*,r_0]$, $\sigma$ a simplex of $\operatorname{lk}_q$, compatibly and without new vertices, send each
vertex $(w,r)$, $r\in\{r_*,r_0\}$, to $rw$, and extend linearly on each simplex. Since the cone parameter is the linear
function equal to $1$ on $\sigma$ on the cone over $\sigma$, this map is a level-preserving piecewise linear homeomorphism of
each prism onto its frustum, and composing the piecewise linear isotopy $(x,r)\mapsto(\gamma_r(x),r)$ with it makes $j'$
piecewise linear.

The new crossings satisfy clause (i). Take the $b-a$ passage, so that $v=\pm(b-a)$ on the passing sheet
(Lemma~\ref{lem:boundaryfield}), and a crossing of the leg with lower support $\{b,c\}$. By Lemma~\ref{lem:monodromy} the
monodromy around this leg is $T(x)=x+\langle\alpha,x\rangle(c-b)$ with $\alpha=u_+-u_-$, up to inversion, which depends on
the orientation of the loop. Since $\langle\alpha,b-a\rangle=0$, $v$ is invariant and extends across the crossing. On
covectors the monodromy is $\xi\mapsto\xi\pm\xi(c-b)\alpha$, and $\ker v$ has lattice
$\mathcal F_v=\{\xi\in M/\Z u_+:\xi(b-a)=0\}$. The map $\mathcal F_v\to\Z$, $\xi\mapsto\xi(c-b)$, is onto, because $b-a$ and
$c-b$ span a saturated sublattice of $u_+^\perp\cap N$, and its kernel is spanned by the class of $\alpha$, which is primitive
in $M/\Z u_+$ because $\{0,u_+,u_-\}$ is a unimodular triangle. In a basis $\{\xi,\alpha\}$ with $\xi(c-b)=1$ the monodromy is
$\bigl(\begin{smallmatrix}1&0\\\pm1&1\end{smallmatrix}\bigr)$, which is conjugate to
$\bigl(\begin{smallmatrix}1&0\\1&1\end{smallmatrix}\bigr)$ in $GL_2(\Z)$ for either sign. The same holds at the leg with lower
support $\{d,a\}$, and for the $c-b$ passage with the roles exchanged. All other clauses of \cite[Def.~7.2]{CBM} are local and
hold where $j'=j$; in $U^{r_0}_q$ the only points of $j'(S)$ on $\Gamma$ are those of $\partial S$ and the new crossings, and
the passing sheet has no interior edges or vertices there. Items (2) and (3) hold for $r_0$ small, since the closures of
the $U^{r}_q$, $r<1$, are pairwise disjoint (Section~\ref{ssec:res-base}), and $\rho$ is read from $\partial S$ and its
boundary field, which are unchanged.
\end{proof}

Clause (i) is read, as in Section~\ref{ssec:cbm}, as conjugacy in $GL_2(\Z)$. With this reading the orientation of the loop
around a leg plays no role in this section, and we fix no convention for it; in particular no statement here depends on
the orientation convention for loops in \cite[Prop.~3.15]{Gross2005}. The shell $U_q\setminus U_q^{r_*}$ may contain crossings, other sheets, interior edges and vertices of $S$, and the
winding of the passing sheet.

\begin{definition}[resolved cycle]\label{def:resolvedcycle}
Let $j'$, $r_*$ be as in Lemma~\ref{lem:NS} and $g$ as in Lemma~\ref{lem:resmove} for this $r_*$. The \emph{resolved
cycle} $\hat S=\hat j(S)$ on $\widehat B$ is defined as follows. Off $\bigcup_qj'^{-1}(U^{r_*}_q)$ put $\hat j:=g\circ j'$, and let $v$ be
transported by $g$. Inside $\widehat U_q$ replace the cone over $\gamma^0_q$ by the \emph{standard half-disc}
$\mathfrak H_q\subset\widehat U_q$: a piecewise linear disc on side $s_q$ of $F_q$ bounded by
$g(\gamma^0_q)\subset\partial\widehat U_q$ and by the path formed by the inner segments of the legs of $\widehat B$ with lower
supports $\{a,b\}$ and $\{c,d\}$ together with $y_0\cup e_q\cup y_\infty$ (for the $c-b$ passage, the legs with lower
supports $\{b,c\}$ and $\{d,a\}$).
\end{definition}

Outside $\bigcup_q\widehat U_q$ the data $(S,\hat j,v)$ satisfy the clauses of \cite[Def.~7.2]{CBM} relative to $\widehat B$,
by Lemma~\ref{lem:resmove}; the image of the shell meets $\widehat\Gamma$ only on the outer segments of the legs at the nodes,
which carry the same monodromy. Inside $\widehat U_q$, $\hat S$ is not a tropical $2$-cycle of $\widehat B$: its boundary turns
at the negative vertices $y_0,y_\infty$ and runs along $e_q$, whose vanishing vector $c-a$ is not $\pm v$. This is the
configuration that produces the exceptional curve. Every crossing point of $\hat S$ lies outside $\bigcup_q\widehat U_q$.

\emph{Normal form near $\partial S$.} The germ of $\hat S$ along $\hat j(\partial S)$ outside $\bigcup_q\widehat U_q$ is
normalised by a piecewise linear isotopy of $\hat j$ that fixes $\hat j(\partial S)$ pointwise, is supported in
a small neighbourhood of $\hat j(\partial S)\setminus\bigcup_q\widehat U_q$ containing no point $\hat j(\mathfrak x_k)$, and
keeps $\hat S$ off $\widehat\Gamma$ outside $\hat j(\partial S)$. Such an isotopy preserves the clauses of
\cite[Def.~7.2]{CBM}, which near $\partial S$ involve only $\hat j(\partial S)$ and the germ of the parallel field $v$ on the
complement of $\widehat\Gamma$. After it:
\begin{enumerate}
\item[(T1)] near every edge of $\partial S$ lying on a leg in a two-face $F_{uu'}$, the sheet of $\hat S$ ending on the edge meets
$F_{uu'}$ only in $\hat j(\partial S)$ and in finitely many arcs, along which the sheet crosses $F_{uu'}$ transversally and which
leave $\hat j(\partial S)$ at distinct interior points of the edge, the \emph{switching points};
\item[(T2)] at every negative vertex $y$, in a two-face $F_{uu'}$, the link of $\hat S$ at $y$ is a tripod lying, apart from its
three ends, in one of the two open hemispheres of the link sphere $\operatorname{lk}_y$ bounded by
$\operatorname{lk}_y\cap F_{uu'}$;
\item[(T3)] at every positive vertex $p$ on $\hat j(\partial S)$, the link of $\hat S$ at $p$ is an arc lying, apart from its
ends, in the open facet that contains the two legs used by $\partial S$.
\end{enumerate}
By (T1), near every point of such an edge other than a switching point the sheet lies on one side of $F_{uu'}$, in $F_u$ or in
$F_{u'}$; in particular it does not run inside $F_{uu'}$ along a stretch of the edge, and it changes side exactly at the
switching points. By Lemma~\ref{lem:NS}, (T2) and (T3), the switching points lie outside $\bigcup_q\widehat U_q$ and away from
the negative and positive vertices.

For (T1): near an edge of $\partial S$ on a leg, away from its end points, both the sheet and $F_{uu'}$ are piecewise linear
surfaces containing the edge, and by general position relative to $\hat j(\partial S)$ \cite[Ch.~5]{RS} a small isotopy fixing
$\hat j(\partial S)$ makes them meet outside the edge in finitely many arcs, along which they cross transversally and which
leave the edge at distinct points. For (T2): the three leg ends lie on the circle $\operatorname{lk}_y\cap F_{uu'}$, and the
link of $\hat S$ is a tripod with these ends, meeting $\widehat\Gamma$ only in them. Tripods in $S^2$ with three fixed ends
form two isotopy classes relative to the ends, distinguished by the cyclic order of the edges at the centre, and each class is
represented in exactly one of the two open hemispheres, since a tripod in a hemisphere has at its centre the cyclic order of
its ends seen from that hemisphere. The level-wise isotopy of the proof of Lemma~\ref{lem:NS}, which here creates no crossings,
moves the germ into that hemisphere. For (T3): the link sphere $\operatorname{lk}_p$ meets the edge of $\nabla^{\check h}$
through $p$ in two points, and the three two-faces through the edge in three arcs joining them, one through each leg end;
these arcs bound three lunes, one in each facet. The link of $\hat S$ joins the two leg ends used by $\partial S$ in the
complement of the third leg end, a disc, so it is isotopic relative to its ends, in that complement, to an arc in the lune
between them, and the level-wise isotopy moves the germ accordingly without creating crossings. This is the normalisation
that \cite[\S7.4]{CBM} makes up to local isotopy.

Finally, choose a closed neighbourhood $\mathcal N_\Gamma$ of $\hat j^{-1}(\widehat\Gamma)=\partial S\cup\{\mathfrak x_k\}$ in
$S$ carrying the local pieces of Section~\ref{ssec:res-local}, and apply a generic piecewise linear isotopy of $\hat j$
relative to $\mathcal N_\Gamma\cup\bigcup_q\hat j^{-1}(\widehat U_q)$, supported off $\widehat\Gamma$, so that off
$\mathcal N_\Gamma$ the surface $\hat S$ avoids the vertices of $\nabla^{\check h}$, meets its edges transversally in finitely
many points (the \emph{edge crossings}) and its two-faces transversally in curves. Such an isotopy preserves the clauses of
\cite[Def.~7.2]{CBM}, which are local near $\widehat\Gamma$.

\subsection{The localized hypersurface}\label{ssec:res-W}
We choose a localization with exact equations near the strata met by the cycle.
Its orbit-preserving isotopy to a holomorphic hypersurface and the topology of its regular regions supply the targets for the chain construction.

Let $\check{\mathcal A}=\Delta^\circ\cap M$ and $\check h_m:=\check h(m)$, so that $\check h_0=0$ and $\check h_m$ is odd and
positive for $m\neq0$. For $R>1$ put
\[
f_R(x)=\sum_{m\in\check{\mathcal A}}(-R)^{-\check h_m}x^m ,
\]
the member of Gross's degeneration $\{\sum_mt^{\check h(m)}x^m=0\}$ at $t=-1/R$, all of whose coefficients are $1$. On the
positive real torus the constant term is $1$ and all other terms are negative. Write
$\operatorname{Log}(x)=-\log|x|/\log R\in N_\R$ for $x$ in the torus $N\otimes\C^*$, and
\[
T_m(x)=\log_R\bigl|(-R)^{-\check h_m}x^m\bigr|=-\check h_m-\langle m,\operatorname{Log}x\rangle .
\]
In this section the tropical variable is $\xi=\operatorname{Log}x\in N_\R$, and the tropical polynomial of the heights
$\check h_m$ is $\xi\mapsto\min_m\bigl(\check h_m+\langle m,\xi\rangle\bigr)=-\max_mT_m$; this is the convention of
Section~\ref{ssec:conventions}, with the monomials in $M$ and the variable in $N_\R$. The constant term dominates exactly over
the interior of $\nabla^{\check h}$, and $\widehat B=\partial\nabla^{\check h}$ is where it ties with other terms.

\begin{definition}[localized hypersurface]\label{def:W}
Let $\beta\colon\R\to[0,1]$ be smooth and non-increasing with $\beta=1$ exactly on $(-\infty,C_1]$ and $\beta=0$ exactly on
$[C_0,\infty)$, where $0<C_1<C_0<\frac14$. Put $\beta_m(x)=\prod_{i\in\check{\mathcal A}}\beta\bigl(T_i(x)-T_m(x)\bigr)$,
\[
\tilde f(x)=\sum_{m\in\check{\mathcal A}}\beta_m(x)(-R)^{-\check h_m}x^m ,\qquad W=\overline{\{\tilde f=0\}}\subset X_{\mathcal F'} ,
\]
and for $s\in[0,1]$ let $V_s$ be the closure of the zero set of
$\sum_m\bigl((1-s)+s\beta_m\bigr)(-R)^{-\check h_m}x^m$. The \emph{active set} at $x$ is
$A(x)=\{m:\beta_m>0\text{ near }x\}$.
\end{definition}

So $V_0=\widehat X_R$, the proper transform of $X_R=\{f_R=0\}$, and $V_1=W$; a monomial is active only where it is within
$C_0$ of the maximum of the $T_i$. The hypersurface $W$ is the localization of $X_R$ in the sense of Mikhalkin
\cite{Mikhalkin}, in the form of \cite[Def.~3.1]{Ya16}: near each point it is the zero set of the monomials that are close to
maximal there, and it is isotopic to $\widehat X_R$ (Lemma~\ref{lem:T0}). In the notation of \cite[\S3]{Ya16}, with
$q_{\rm Y}=-R$, $k_m=q_{\rm Y}^{v_m}$ and $v_m=-\check h_m$, the function $f_R$ is $f_{q_{\rm Y}}$, $\tilde f$ is
$\tilde f_{q_{\rm Y}}$ and $W$ is $W_{q_{\rm Y}}$. The leading coefficients $c_m$ of loc.\ cit.\ are all $1$, so the correction
term in Yamamoto's interpolating family $\tilde f_{q_{\rm Y},s}$ vanishes, and $V_s$ is the member $V_{q_{\rm Y},s'}$ of that
family for every $s'$ with $d(s')=s$ \cite[proof of Thm.~3.9]{Ya16}.

\begin{lemma}\label{lem:Ufan}
\begin{enumerate}
\item For $C_0$ small, every $V_s$ meets only torus orbits $O(\sigma)$ of $X_{\mathcal F'}$ with $\sigma$ contained in the cone
over a face of $\Delta$ of dimension at most two.
\item Every cone of $\mathcal F'$ contained in the cone over a face of $\Delta$ of dimension at most two is unimodular; so
$X_{\mathcal F'}$ is smooth along the orbits met by the $V_s$.
\item For $C_0$ small and $R$ large, for every $s\in[0,1]$ and every torus orbit $O(\sigma)$ met by $V_s$, the restriction to
$O(\sigma)$ of the defining function of $V_s$ has $0$ as a regular value, and its derivatives along the radial and angular
fundamental vector fields of $N\otimes\C^*$ span $\C$ at every point of $V_s\cap O(\sigma)$. Hence every $V_s$ is a smooth
hypersurface transverse to every torus orbit it meets. Near each orbit the active monomials at a point of $W$ are the vertices
of one simplex of the triangulation of $\Delta^\circ$ induced by the heights, and the equation of $W$ is the localized
hyperplane \cite[Def.~3.2]{Ya16} in the ratios of the active monomials, independent of the remaining coordinates
\cite[Thm.~3.9(2),(3)]{Ya16}.
\end{enumerate}
\end{lemma}
\begin{proof}
(1) In Cox coordinates $(Z_\varrho)$, indexed by the rays $\varrho$ of $\mathcal F'$ with primitive generators $n_\varrho$, the
defining function is $\sum_m\beta^{(s)}_m(-R)^{-\check h_m}\prod_\varrho Z_\varrho^{\langle m,n_\varrho\rangle+1}$ with
$\beta^{(s)}_m=(1-s)+s\beta_m$, and on $O(\sigma)$ only the terms with $m$ in the dual face
$\sigma^*=\{m\in\Delta^\circ:\langle m,n_\varrho\rangle=-1,\ \varrho\in\sigma(1)\}$ survive. If $\sigma$ is not contained in the
cone over a face of dimension at most two, $\sigma^*$ is a vertex $u$ of $\Delta^\circ$. Approaching $O(\sigma)$,
$\operatorname{Log}$ tends to infinity in a direction $n$ of the relative interior of $\sigma$, and $T_m-T_u\to-\infty$ for
$m\ne u$ because $\langle m-u,n\rangle>0$; so $\beta_u=1$ near $O(\sigma)$ and the restriction is a nonvanishing monomial.
(2) Such a cone is the cone over a primitive segment or a unimodular triangle at lattice height one (Lemma~\ref{lem:regref}),
hence unimodular. (3) The last sentence is \cite[Thm.~3.9(2),(3)]{Ya16}. Yamamoto assumes that all coefficients are nonzero,
that the heights are integral, that the subdivision of $\Delta^\circ$ induced by the heights is unimodular, which holds because
$\mathcal T$ is unimodular (Theorem~\ref{thm:U}), and that the refinement of the normal fan of $\Delta^\circ$ is unimodular;
statements (2) and (3) of his theorem are local near each orbit met by $W$, and by (1) and (2) above the charts
$U_\sigma\cong\C^k\times(\C^*)^{4-k}$ of the relevant cones are those of his setting. The regularity statement is asserted in
the proof of \cite[Thm.~3.9]{Ya16} as a consequence of $\det\mathbf J\neq0$ for a $2\times2$ matrix $\mathbf J$, ``by the
concrete calculation'', which is not written there. We do not rely on that assertion: the calculation is carried out in
Appendix~\ref{app:Ufan}, which proves (3), and $\det\mathbf J>0$ there for $R$ large, uniformly in $s$ and in the point.
\end{proof}

\begin{lemma}[exact local symmetry]\label{lem:symm}
Let $x$ be a point of $W$ in the dense torus. The subgroup $\bigl((A(x)-A(x))^\perp\cap N\bigr)\otimes U(1)$ of the compact
torus $N\otimes U(1)$ preserves $W$ near $x$ and near the points of $W$ on the toric strata in the closure of a neighbourhood
of $x$.
\end{lemma}
\begin{proof}
Each $\beta_m$ depends only on the moduli $|x^i|$, so for $\lambda$ in the compact torus
$\tilde f(\lambda x)=\sum_m\beta_m(x)(-R)^{-\check h_m}\lambda^mx^m$. If $\lambda^m$ is constant on $A(x)$, then
$\tilde f(\lambda x)=\lambda^{m_0}\tilde f(x)$ near $x$, for any $m_0\in A(x)$.
\end{proof}

\begin{lemma}[no positive points on the strata]\label{lem:stratum}
The closure of $W\cap(\R_{>0})^4$ in $X_{\mathcal F'}$ meets no toric stratum.
\end{lemma}
\begin{proof}
On $O(\sigma)$, $\sigma\neq0$, the constant term vanishes, since its Cox monomial is $\prod_\varrho Z_\varrho$. At a
nonnegative real point of $O(\sigma)$ every surviving term, $m\in\sigma^*\setminus\{0\}$, is a non-positive real number, and
the largest active one has $\beta_m=1$; so the restriction is strictly negative.
\end{proof}

A chain that ends on a toric stratum and starts on the positive real points must therefore leave the positive reals. This is
why the local pieces of Section~\ref{ssec:res-local} contain collar paths, paths in $W$ from a point of the positive section
of Lemma~\ref{lem:pos} to the critical locus of a local model.

\begin{lemma}[the positive section]\label{lem:pos}
For $R$ large, $W_{>0}=W\cap(\R_{>0})^4$ is a smooth hypersurface of $(\R_{>0})^4$ along which $W$ is smooth, and
$\operatorname{Log}$ maps it homeomorphically onto a star-shaped sphere $\Sigma_{>0}$ lying within radial distance
$2C_0t_{\max}/(1-C_0)$ of $\widehat B$, where $t_{\max}=1+\max_{y\in\nabla^{\check h}}|y|$. Radial projection identifies
$\Sigma_{>0}$ with $\widehat B$. The inverse $\mathfrak s\colon\widehat B\to W$ of the composite is continuous, and
$\operatorname{Log}\mathfrak s(b)$ lies on the ray through $b$.
\end{lemma}
The proof is in Appendix~\ref{app:pos}; it uses only $\check h_m\ge1$ for $m\neq0$, the parity of $\check h_m$, that all
coefficients of $f_R$ have modulus one, and $R$ large. We call $\mathfrak s$ the \emph{positive section}.

\emph{Charts of the faces.} Let $\Phi$ be a facet, a two-face or an edge of $\nabla^{\check h}$, and
$A_\Phi=\{0\}\cup\{u:\Phi\subset F_u\}$ its \emph{tie set}, of type (a) $\{0,u\}$, (b) $\{0,u,u'\}$ or (c)
$\{0,u_1,u_2,u_3\}$. For types (b), (c) a \emph{fan point} of $\Phi$ is a vertex $n$ of $\mathcal F'$ whose vertex region
$\widetilde R_n$ meets $\Phi$; it satisfies $\langle a,n\rangle=-1$ for $a\in A_\Phi\setminus0$. Since $\mathcal T$ is
unimodular, $A_\Phi\setminus0$ extends to a basis of $M$, so $N=A_\Phi^\perp\oplus\Z n\oplus C$ for a lattice complement
$C$. A \emph{chart} of $(\Phi,n)$ is a linear isomorphism
$(\pi_\parallel,\pi_\perp)\colon N_\R\to(A_\Phi^\perp)_\R\times\R^{A_\Phi\setminus0}$ with
\[
\pi_\perp(y)=\bigl(\langle a,y\rangle+\check h_a\bigr)_{a\in A_\Phi\setminus0},\qquad \pi_\parallel(n)=0 ,
\]
and $\pi_\parallel$ a projection onto $(A_\Phi^\perp)_\R$; the condition $\pi_\parallel(n)=0$ is imposed in types (b) and
(c), and in type (a) $\pi_\parallel$ is the projection along a vector $n_F$ with $\langle u,n_F\rangle=-1$. Near $\Phi$ the
polytope is $\{\pi_\perp\ge0\}$, so in the chart $\nabla^{\check h}=(A^\perp_\Phi)_\R\times\mathcal O_\perp$ and
$\widehat B=(A^\perp_\Phi)_\R\times\partial\mathcal O_\perp$, with $\mathcal O_\perp=\R^{A_\Phi\setminus0}_{\ge0}$;
$\partial\mathcal O_\perp$ is a point, two half-lines or three quadrants. The \emph{tentacle} from $y$ is $y+\R_{\ge0}n$; in the
chart it is $\{\pi_\parallel(y)\}\times(\pi_\perp(y)+\R_{\ge0}(-1,\dots,-1))$, and it runs into the toric divisor $D_n$.
Yamamoto's tropical contraction $\varpi$ \cite{Ya21} maps the tropical hypersurface onto $\widehat B$ by the identity on the
facet interiors and along the tentacles elsewhere.

\begin{definition}[product-shaped sets and regions]\label{def:regions}
An open set $U\subset\widehat B$ is \emph{product-shaped} at $\Phi$ if in a chart of $(\Phi,n)$ it is
$\{\pi_\parallel\in U_\parallel,\ \pi_\perp\in U_\perp\}$ with $U_\parallel$ convex and
$U_\perp=\partial\mathcal O_\perp\cap\{|\pi_\perp|_\infty<r\}$, and its distance to every face not containing $\Phi$ and to
$\widehat\Gamma$ exceeds $\ell_0$. For such $U$ let $\Omega_U=\varpi^{-1}(U)^\eta$, the $\eta$-neighbourhood in the sup-norm of
the chart; in the chart
\[
\Omega_U=\{\pi_\parallel\in U_\parallel^\eta,\ \pi_\perp\in\Omega_\perp\},\qquad
\Omega_\perp=S_{d,r}^{\eta},\qquad d=|A_\Phi\setminus0|,
\]
where, for $p=(p_1,\ldots,p_d)$ and $\mu(p)=\min(0,p_1,\ldots,p_d)$, we set
\[
 S_{d,r}=\{p:\text{the minimum defining }\mu(p)\text{ is attained at least twice},\quad
                    p_i-\mu(p)<r\ (1\le i\le d)\}.
\]
Thus $S_{1,r}=\{0\}$; $S_{2,r}$ consists of two truncated positive rays and the negative diagonal ray;
and $S_{3,r}$ consists of the truncated boundary of the positive orthant together with its tentacles from the rays and
vertex. The contraction is $p\mapsto p-\mu(p)(1,\ldots,1)$.
The \emph{region} of
$U$ is $W_U=W\cap\operatorname{Log}^{-1}(\Omega_U)$ together with its limit points on the divisor $D_n$ at the ends of the
tentacles. Here $\eta>C_0$ and $\ell_0=C_{\rm Lip}(C_{\rm ch}\eta+C_{\rm act}C_0)$, where $C_{\rm ch}$ bounds the distortion
between the chart norms and the Euclidean norm, $C_{\rm Lip}$ is the Lipschitz constant of $\varpi$ over pairs of cells, and
$C_{\rm act}$ is such that active monomials lie within $C_{\rm act}C_0$ of the maximum.
\end{definition}

The regions play, for $W$, the role that the preimages $U\times T^3$ of small open sets $U$ of the smooth part of the base
play for a torus fibration, as in the topological models of \cite{CBM}; Lemma~\ref{lem:regions} makes this precise at the
level of homotopy.

\begin{lemma}[regions]\label{lem:regions}
For $R$ large and every product-shaped $U$ of type (a), (b) or (c):
\begin{enumerate}
\item $W_U\cong(A_\Phi^\perp\otimes U(1))\times U_\parallel^\eta\times L_U$, where $L_U$ is a point in type (a), homotopy
equivalent to a circle in type (b), and homotopy equivalent to a two-torus in type (c). Hence $W_U$ is aspherical, and
$\pi_1(W_U)\cong\Lambda_U$, the lattice $u^\perp\cap N$ in type (a) and $N/\Z n$ in types (b) and (c).
\item If $U'$ is product-shaped and $W_{U'}\subset W_U$, the inclusion induces on $\pi_1$ the parallel transport of $\Lambda$
from $U'$ to $U$.
\item The same holds for the sheared charts $(\pi_\parallel+\lambda\circ(\pi_\perp-\pi_\perp(0)),\pi_\perp)$ with $\lambda$
linear and $\lambda(\pi_\perp(n)-\pi_\perp(0))=0$.
\end{enumerate}
\end{lemma}
The proof is in Appendix~\ref{app:regions}. In type (b), $L_U$ is a truncated localized line capped at the divisor $D_n$; in
type (c), it is a truncated localized plane, $\mathbb P^2$ minus four lines, capped at $D_n$. By Lemma~\ref{lem:pos},
$\mathfrak s(U)\subset W_U$ once $\eta$ exceeds $C_{\rm ch}$ times the radial width of Lemma~\ref{lem:pos}, and loops
$\mathfrak s(\gamma)$ with $\gamma$ in a contractible $U$ are contractible in $W_U$.

\subsection{The domain}\label{ssec:res-domain}
We construct the three-dimensional domain by attaching circles, triangles and tetrahedra over the cells of the tropical cycle.
Balancing cancels the boundaries along interior strata; the remaining boundary components lie over the nodes.

Let $\mathfrak T$ be the triangulation of $S\setminus\operatorname{int}\mathcal N_\Gamma$ of Lemma~\ref{lem:AC} below. The
cell complex $\mathfrak D(S)$ is glued from pieces over the cells of $\mathfrak T$ and from local domains over $\mathcal N_\Gamma$.

\begin{definition}[the domain]\label{def:domain}
The \emph{domain} $\mathfrak D=\mathfrak D(S)$ is the CW complex obtained by gluing the following pieces.
\begin{itemize}
\item Over each cell $c$ of $\mathfrak T$ in a sheet, $c\times S^1$, the circle being $\R v/\Z v$, oriented by $v$.
\item Over each cell $c$ of $\mathfrak T$ on an interior edge, with fields $v_1,v_2,v_3$ on the adjacent sheets and
$s_1v_1+s_2v_2+s_3v_3=0$ (clause (g)): $c\times\mathcal Q$, where $\mathcal Q$ is a $2$-simplex with its three vertices identified
to one point, whose three edges are glued to the circles $G_k=\R v_k/\Z v_k$ of the three sheets, and which is oriented so that
$\partial\mathcal Q=s_1G_1+s_2G_2+s_3G_3$ with $G_k$ oriented by $v_k$. It is the domain of the flat triangle of
\cite[Lemma~7.7]{CBM}: the linear map sending the edges of a $2$-simplex to $s_kv_k$, followed by the quotient by the lattice.
\item Over an interior vertex, where four interior edges and six sheets meet: a $3$-simplex whose four faces are glued to the
triangles $\mathcal Q$ of the four edges. This is possible because the six fields are the edge vectors, and the four balancing
relations the face relations, of the image of a $3$-simplex under a linear map, which is nondegenerate by clause (j)
\cite[\S7.7, Step~4]{CBM}.
\item Over $\mathcal N_\Gamma$, the local domains of Section~\ref{ssec:res-local}: a disc times an arc along each leg of
$\partial S$ (over the leg pieces and the switching discs), a cone over a triangle at each boundary vertex, a copy of
$\C\times\R$ at each positive vertex met by $\partial S$, a solid torus at each crossing $\mathfrak x_k$, and, at each node, the
\emph{node domain} $(\mathfrak H^\circ_q\times S^1)/\!\sim$, where $\mathfrak H^\circ_q\subset\mathfrak H_q$ is the part of
$\hat j(\mathcal N_\Gamma)$ in the standard half-disc, a collar of its boundary path, and the circle is collapsed over the two
leg segments.
\end{itemize}
Each piece is oriented as (sheet)$\times$(circle oriented by $v$), and $[\mathfrak D]$ denotes the sum of the oriented top
cells. For a node $q$ on $\partial S$ let $\partial_q\mathfrak D=(e_q\times S^1)/\!\sim$, the circle collapsed at $y_0$ and
$y_\infty$, a $2$-sphere. The frontier of $\mathcal N_\Gamma$ in $S$ is the \emph{inner boundary}.
\end{definition}

Over the smooth part of $S$ away from $\mathcal N_\Gamma$, $\mathfrak D(S)$ is the circle bundle with fibre $\R v/\Z v$ of
\cite[\S7]{CBM}, and over the interior edges and vertices it is their completion by flat triangles and tetrahedra. Each local
domain meets the cells over the regular part $S\setminus\operatorname{int}\mathcal N_\Gamma$ over its part of the inner boundary.

\begin{lemma}[the domain]\label{lem:domain}
$\partial[\mathfrak D]=\sum_q[\partial_q\mathfrak D]$, the sum over the nodes $q$ with $p_q\in j(\partial S)$, where $e_q$
carries the boundary orientation of $\mathfrak H_q$.
\end{lemma}
\begin{proof}
Interior faces cancel in pairs if the local orientations agree across them. Over the sheets this holds because the
orientations are products. Along an interior edge, let $c$ be a one-cell of $\mathfrak T$ on the edge, oriented as the edge.
Then $\partial(c\times\mathcal Q)=\partial c\times\mathcal Q-c\times\partial\mathcal Q=\partial c\times\mathcal Q-\sum_ks_k\,c\times
G_k$. The term $\partial c\times\mathcal Q$ cancels between adjacent cells of the edge, and $s_k\,c\times G_k$ is the
contribution of $c$ to the boundary of the cylinders over the $k$-th sheet, because $s_k$ is by definition the sign with which
$c$ occurs in the boundary of the $k$-th sheet; so the three cylinders are cancelled. At an interior vertex the boundary of the
$3$-simplex is the sum of the four triangles with the signs of the balancing conditions. At the local domains the pieces are
glued along the inner boundaries of Section~\ref{ssec:res-local}, and over the legs, boundary vertices, positive vertices and
crossings the circle collapses or closes up without creating boundary. What remains is the part of
$\partial\mathfrak H^\circ_q\times S^1$ over $e_q$, where the circle $\R v/\Z v$ does not collapse, because the vanishing vector
of $e_q$ is $c-a\neq\pm v$.
\end{proof}

\subsection{The local pieces}\label{ssec:res-local}
Each local domain is mapped into $W$ by an explicit map, built in the exact local models of Lemma~\ref{lem:Ufan}(3), whose
restriction to the inner boundary with the regular part consists of orbits of the circle $S^1_v=\R v/\Z v\subset N\otimes U(1)$
through points of the positive section, over base points in the sets $V_c$ of Lemma~\ref{lem:AC}. Since the positive points
miss the strata (Lemma~\ref{lem:stratum}), each piece contains collar paths from a point of $\mathfrak s$ to the critical locus
of the local model. For a path $\gamma$ in $W$ and a character $x^m$ that does not vanish on $\gamma$, a \emph{lifted phase} of
$x^m$ along $\gamma$ is a continuous function $\theta$ with $x^m(\gamma(\cdot))\in\R_{>0}e^{i\theta(\cdot)}$; its total change
along $\gamma$ does not depend on the choice. If $x^m$ vanishes only at the end point of $\gamma$ and its phase is constant near
that point, we use the limit.

\emph{Adapted coordinates.} For each two-face $F_{uu'}$ of $\nabla^{\check h}$ we fix an ordering $(u,u')$ and a cocharacter
$\nu=\nu_{uu'}\in N$ with
\[
\langle u,\nu\rangle=0,\qquad\langle u',\nu\rangle=1;
\]
it exists because $\{0,u,u'\}$ is a unimodular triangle, so that $u,u'$ extend to a basis of $M$. Let a leg in $F_{uu'}$ have
lower support $\{n,n'\}$, and $\sigma=\operatorname{cone}(n,n')\in\mathcal F'$. Choose $m_1\in M$ with $\langle m_1,n\rangle=1$,
$\langle m_1,n'\rangle=0$ and $\langle m_1,\nu\rangle=0$ (adjust any solution of the first two conditions by a multiple of
$u'-u\in\sigma^\perp$), put $m_2=-u-m_1$, so that $\langle m_2,n\rangle=0$, $\langle m_2,n'\rangle=1$ and
$\langle m_2,\nu\rangle=0$, and let $\kappa$ generate $\sigma^\perp\cap\nu^\perp\cap M$, which has rank one. Then
$(m_1,m_2,u'-u,\kappa)$ is a basis of $M$ with dual basis $(n,n',\nu,\cdot)$, and
\[
z_1=x^{m_1},\quad z_2=x^{m_2},\quad\mathfrak t=x^{u'-u},\quad w=x^\kappa,\qquad z_1z_2=x^{-u},
\]
are coordinates on $U_\sigma\cong\C^2\times(\C^*)^2$. Near the leg only $0,u,u'$ are active, and multiplying $\tilde f$ by
$x^{-u}$ shows that $W\cap U_\sigma$ is
\begin{equation}\label{eq:legmodelA}
\beta_0\,z_1z_2=\alpha_u+\alpha_{u'}\mathfrak t,\qquad\alpha_u=\beta_uR^{-\check h_u},\quad\alpha_{u'}=\beta_{u'}R^{-\check h_{u'}},
\end{equation}
where $\beta_0,\alpha_u,\alpha_{u'}\ge0$ are functions of $|x|$ (the signs because $\check h_u,\check h_{u'}$ are odd), and
$\beta_0$ vanishes where the constant term is cut off, in particular near $D_n\cup D_{n'}$. The circle $S^1_v$, $v=n'-n$, acts
with weights $(-1,1,0,0)$ on $(z_1,z_2,\mathfrak t,w)$, exactly by Lemma~\ref{lem:symm}, and freely off
$\mathrm{Crit}=\{z_1=z_2=0,\ \mathfrak t=\mathfrak t_0\}$, $\mathfrak t_0=-\alpha_u/\alpha_{u'}<0$. We call $\mathrm{Crit}$ the
\emph{critical locus} of the leg model: it is the fixed locus of $S^1_v$ in $W\cap U_\sigma$, the curve over which the circle of
the boundary field collapses.

\emph{Sides and chambers.} By (T1), near a point $e$ of an edge of $\partial S$ on the leg that is not a switching point, the
sheet of $\hat S$ ending on the leg lies on the side of $F_u$ or of $F_{u'}$; its \emph{side} at $e$ is $s=+1$ or
$-1$ accordingly. At every point of $W$ near the leg, $\alpha_u-\alpha_{u'}|\mathfrak t|$ has the sign of
$T_u-T_{u'}=\check h_{u'}-\check h_u-\log_R|\mathfrak t|$, because
$\alpha_{u'}|\mathfrak t|/\alpha_u=R^{T_{u'}-T_u}\beta_{u'}/\beta_u$ and $\beta_{u'}\le\beta_u$ if and only if $T_{u'}\le T_u$.
The \emph{chamber} of side $s$ at the leg is
\[
\mathrm{Ch}_s=W\cap U_\sigma\cap\{z_1z_2\in s\R_{\ge0},\ \mathfrak t\in\R_{<0},\ w\in\R_{>0}\}
=W\cap U_\sigma\cap\{x^{-u}\in s\R_{\ge0},\ \mathfrak t\in\R_{<0},\ w\in\R_{>0}\};
\]
it is $S^1_v$-invariant, contains $\mathrm{Crit}\cap\{w>0\}$, and on it
$s\beta_0|z_1z_2|=\alpha_u-\alpha_{u'}|\mathfrak t|$.

\emph{Coefficients near a leg.} Put $\mu=T_0-T_u$ and $\vartheta=T_u-T_{u'}$. Near the leg these depend only on
$|z_1z_2|=R^{\mu-\check h_u}$ and on $|\mathfrak t|=R^{\check h_{u'}-\check h_u-\vartheta}$; $\vartheta=0$ exactly where
$|\mathfrak t|=|\mathfrak t_0|$, and the points of $\widehat B$ near the leg on side $s$ have $s\vartheta>0$.
Since only $0$, $u$ and $u'$ are active near the leg,
\begin{equation}\label{eq:legcoeff}
\beta_0=\beta(-\mu)\,\beta(-\mu-\vartheta),\qquad R^{\check h_u}\alpha_u=\beta(\mu)\,\beta(-\vartheta),\qquad
R^{\check h_u}\alpha_{u'}|\mathfrak t|=R^{-\vartheta}\,\beta(\mu+\vartheta)\,\beta(\vartheta).
\end{equation}

\begin{lemma}[real points of a chamber]\label{lem:chamber}
For $\vartheta'>0$ put
\[
f_{\vartheta'}(\xi)=\beta(-\xi)\,\beta(-\xi-\vartheta')\,R^{\xi}-\beta(\xi)+R^{-\vartheta'}\beta(\vartheta')\,\beta(\xi+\vartheta') .
\]
For $R$ large, every $f_{\vartheta'}$ has exactly one zero $X(\vartheta')$; $X(\vartheta')\in(-C_0,C_0)$, $X$ is continuous, and
$X(\vartheta')\to-C_0$ as $\vartheta'\to0$. Near the leg, a point of $\mathrm{Ch}_s$ with $z_1z_2\neq0$, $z_2>0$ and
given $w$ and $|z_1|/|z_2|$ is determined by $(\mu,\vartheta)$, and such points are:
\begin{enumerate}
\item for $s\vartheta>0$, the points with $T_0-T_{u_s}=X(s\vartheta)$, where $u_{+1}=u$ and
$u_{-1}=u'$;
\item for $\vartheta=0$, that is $\mathfrak t=\mathfrak t_0$, the points with $|z_1z_2|\le R^{-C_0-\check h_u}$.
\end{enumerate}
\end{lemma}
\begin{proof}
On $\mathrm{Ch}_s$ the equation is $s\beta_0|z_1z_2|=\alpha_u-\alpha_{u'}|\mathfrak t|$. For $s=1$
and $\vartheta\ge0$ we have $\beta(-\vartheta)=1$, and by \eqref{eq:legcoeff}, after multiplication by $R^{\check h_u}$, the
equation reads $f_\vartheta(\mu)=0$. For $s=-1$ and $\vartheta=-\vartheta'\le0$ we have $\beta(\vartheta)=1$, and after
multiplication by $-R^{\check h_u-\vartheta'}$ it reads $f_{\vartheta'}(\mu')=0$ with $\mu'=\mu-\vartheta'=T_0-T_{u'}$. At
$\vartheta=0$ it reads $\beta(-\mu)^2R^{\mu}=0$, which holds exactly when $\mu\le-C_0$. So (1) and (2) follow from the
statements on $f_{\vartheta'}$.

Let $\vartheta'>0$. For $\xi\le-C_0$, $\beta(-\xi)=0$ and $\beta(\xi)=1$, so $f_{\vartheta'}(\xi)\le
R^{-\vartheta'}\beta(\vartheta')-1<0$; for $\xi\ge C_0$, $f_{\vartheta'}(\xi)=R^\xi>0$. For $\xi>-C_0$ we have
$\beta(-\xi-\vartheta')\ge\beta(-\xi)>0$, and since $\beta'\le0$,
\[
f'_{\vartheta'}(\xi)\ \ge\ R^\xi\log R\cdot\beta(-\xi)^2+R^{-\vartheta'}\beta(\vartheta')\,\beta'(\xi+\vartheta').
\]
The last term vanishes unless $C_1<\xi+\vartheta'<C_0$ and $\vartheta'<C_0$. Then $\xi>C_1-C_0$, so
$\beta(-\xi)\ge\beta(C_0-C_1)>0$ and $R^\xi>R^{C_1-\vartheta'}$, and
$f'_{\vartheta'}(\xi)\ge R^{-\vartheta'}\bigl(R^{C_1}\log R\cdot\beta(C_0-C_1)^2-\max|\beta'|\bigr)$, which is positive for $R$
large. So $f_{\vartheta'}$ has positive derivative on $(-C_0,\infty)$ and exactly one zero $X(\vartheta')\in(-C_0,C_0)$, which
depends smoothly on $\vartheta'$ by the implicit function theorem. As $\vartheta'\to0$, $f_{\vartheta'}$ converges uniformly on
$[-C_0,C_0]$ to $\beta(-\xi)^2R^\xi$, which is bounded below by a positive constant on $[-C_0+\eta',C_0]$ for every $\eta'>0$;
so $X(\vartheta')\to-C_0$.
\end{proof}

\begin{definition}[leg collar paths]\label{def:collarpath}
Let $e$ be a point of an edge of $\partial S$ on the leg that is not a switching point, and $s$ its side. Over the
transverse segment $[0,\varsigma_0]$ of $\hat S$ at $e$, the \emph{leg collar path} $\gamma_e(\varsigma)$ is the path in $W$
that
\begin{enumerate}
\item[(i)] for $\varsigma\in[\varsigma_0/2,\varsigma_0]$ follows $\mathfrak s$;
\item[(ii)] then moves along the positive real points of $W$, with $|z_1|/|z_2|$ and $w$ fixed, changing $|\mathfrak t|$
monotonically until $s\vartheta\ge C_0$ (it is constant if this holds at its start);
\item[(iii)] then turns $\arg\mathfrak t$ from $0$ to $\pi$ with $|\mathfrak t|$, $|z_1|$, $z_2$ and $w$ fixed;
\item[(iv)] then moves inside $\mathrm{Ch}_s$ with $z_2>0$ and with $w$ and $|z_1|/|z_2|$ fixed: first through the
points of Lemma~\ref{lem:chamber}(1), $s\vartheta$ decreasing to $0$, then through those of
Lemma~\ref{lem:chamber}(2), $|z_1z_2|$ decreasing from $R^{-C_0-\check h_u}$ to $0$. It ends at the \emph{point of
$\mathrm{Crit}$ over $e$}, the point of $\mathrm{Crit}$ (where $z=0$ and $\mathfrak t=\mathfrak t_0$) with the value
$w=w(\mathfrak s(e,\varsigma_0/2))$ fixed along the collar path.
\end{enumerate}
\end{definition}

Over a transverse segment, the collar paths play the role of the sections $\check\sigma$ with $\check\sigma(\partial S)$ in the
critical locus that Casta\~no-Bernard and Matessi use in their local models \cite[\S7.4, \S7.5]{CBM}, and the leg piece below is
the orbit of such a section under $S^1_v$. The end point depends continuously on $e$ because $\mathfrak s$ does. In (ii) the positive real point with given
$|\mathfrak t|$, $w$ and $|z_1|/|z_2|$ is unique and depends continuously on $|\mathfrak t|$: by \eqref{eq:legcoeff} its
equation is $\beta_0R^\mu=R^{\check h_u}(\alpha_u+\alpha_{u'}|\mathfrak t|)$, whose left side vanishes for
$\mu\le\mu^*:=\max(-C_0,-C_0-\vartheta)$ (for $\vartheta<0$ the factor $\beta_0$ also vanishes on $(-C_0,-C_0-\vartheta]$) and
has positive derivative for $\mu>\mu^*$, while the right side is positive for $\mu\le\mu^*$, non-increasing in $\mu$, and zero
for $\mu$ large. At the end of (ii), $\beta(s\vartheta)=0$, so by \eqref{eq:legcoeff} $\alpha_{u'}=0$ if
$s=1$ and $\alpha_u=0$ if $s=-1$. The coefficients of \eqref{eq:legmodelA} depend only on $|z_1z_2|$ and
$|\mathfrak t|$, so they are constant along (iii), and (iii) is explicit: for $s=1$, $z_1=\alpha_u/(\beta_0z_2)$ does
not move, and for $s=-1$, $z_1=\alpha_{u'}\mathfrak t/(\beta_0z_2)$ turns with $\mathfrak t$ from the positive to the
negative real axis ($\beta_0>0$ at these points, since the right side of the equation is positive at the end of (ii)). Since
$\beta(s\vartheta)=0$ there, the equation of the positive real points at the end of (ii) coincides with that of
$\mathrm{Ch}_s$, so the end point of (iii) is the point of Lemma~\ref{lem:chamber}(1) with this value of
$s\vartheta$, where (iv) starts. By Lemma~\ref{lem:chamber}, (iv) is a continuous path, and (i)--(iv) depend
continuously on $e$ along each open interval of the edge between switching points. Along the whole collar path the lifted phase
of $\mathfrak t$ increases by $\pi$, those of $z_2$ and $w$ do not change, and that of $z_1$ increases by $(1-s)\pi/2$
(in (iv), $z_1$ and $z_2$ keep their phases until they vanish at the end point); hence the lifted phase of a character $x^m$
changes by
\begin{equation}\label{eq:collarphase}
\pi\Bigl(\langle m,\nu\rangle+\tfrac{1-s}{2}\langle m,n\rangle\Bigr),
\end{equation}
and $x^{-u}=z_1z_2$ ends in $s\R_{\ge0}$. The \emph{leg piece} is $S^1_v\cdot\{\gamma_e(\varsigma)\}$, a disc times an
arc: the circle collapses at $\varsigma=0$, on $\mathrm{Crit}$, so the piece has no boundary over the leg. Its outer part
consists of $S^1_v$-orbits through points of $\mathfrak s$.

\emph{Switching discs.} At a switching point $e^*$ the side changes, and the leg collar paths of the two sides end (iii) with
$z_1$ of opposite signs, so they do not form a continuous family across $e^*$. For each switching point $e^*$ we choose points
$e'$ and $e''$ of the edge on either side of $e^*$, with no other switching point between them, and call \emph{switching disc}
the closed disc $\mathbb B_{e^*}\subset S$ bounded by the segment $[e',e'']$ of $\partial S$, the transverse segments at $e'$ and
$e''$, and an arc of the inner boundary. The \emph{switching piece} is $S^1_v\cdot\Theta_{e^*}$ over $\mathbb B_{e^*}$, where
$\Theta_{e^*}\colon\mathbb B_{e^*}\to W$ is given by Lemma~\ref{lem:collar}(4): it agrees with $\gamma_{e'}$ and $\gamma_{e''}$
on the transverse segments and with $\mathfrak s$ on the inner boundary arc, and maps each point $e$ of $[e',e'']$ to the point of
$\mathrm{Crit}$ over $e$. Like a leg piece it is a disc times an arc; the leg pieces are taken over the edge with the interiors
of the switching discs removed.

\emph{Boundary vertices of $S$} (negative vertices). Let the lower support be the unimodular triangle $\{n_1,n_2,n_3\}$ in
$F_{uu'}$, and $\sigma=\operatorname{cone}(n_1,n_2,n_3)$. Let $Z_j=x^{m^{(j)}}$ with $\langle m^{(j)},n_k\rangle=\delta_{jk}$ and
$\langle m^{(j)},\nu\rangle=0$; then $Z_1Z_2Z_3=x^{-u}$, $(Z_1,Z_2,Z_3,\mathfrak t)$ are coordinates on
$U_\sigma\cong\C^3\times\C^*$, and $W$ is $\beta_0Z_1Z_2Z_3=\alpha_u+\alpha_{u'}\mathfrak t$. On the leg with lower support
$\{n_k,n_l\}$ the adapted coordinates may be taken to be $z_1=Z_k$, $z_2=Z_lZ_j$, $w=Z_j$, $\{j,k,l\}=\{1,2,3\}$. The torus
$T^2=(u,u')^\perp\otimes U(1)=\{\lambda_1\lambda_2\lambda_3=1\}$ acts exactly; the three legs are the coordinate axes $L_j$,
whose stabilisers $G_j$ are the circles of the three boundary fields. By (T2) the three sheets at the vertex have one side
$s$. The \emph{vertex path} is a path in $W$ over the segment of the interior edge of $S$ leaving the vertex, from a
point of $\mathfrak s$ to the point $Z=0$, $\mathfrak t=\mathfrak t_0$ of $\mathrm{Crit}$ over the vertex, which turns
$\arg\mathfrak t$ from $0$ to $\pi$ as in (iii) and ends with $|Z_1|=|Z_2|=|Z_3|\to0$ and $\arg Z_j=0$ (side $+1$) or $\pi/3$
(side $-1$), so that $Z_1Z_2Z_3=x^{-u}$ ends in $s\R_{\ge0}$. Let $\mathbb B_j$ be the corner of the $j$-th sheet at the
vertex: a closed disc in the sheet bounded by a segment of its leg ending at the vertex, a segment of the interior edge, a
transverse segment of the sheet at the other end of the leg segment, and an arc of the inner boundary; the leg segment contains no
switching point. The piece is the union of the leg pieces over the rest of the legs, the three sheet pieces $G_j\cdot\Theta_j$
over the $\mathbb B_j$, with $\Theta_j$ given by Lemma~\ref{lem:collar}(2), and the cone, over the vertex path, on the triangle
$\mathcal Q\subset T^2$ with $\partial\mathcal Q=G_1+G_2+G_3$ \cite[\S7.5]{CBM}. The two-cells $\mathcal Q\cdot\mathfrak s(b)$
over the interior edge leaving the vertex lie in $T^2\cdot\mathfrak s(b)$.

\emph{Positive vertices.} Let $\partial S$ turn at the positive vertex $p$ on the edge of $\nabla^{\check h}$ dual to
$\{u_1,u_2,u_3\}$, with lower support $\{n,n'\}$, along the legs $\Gamma_{12}\subset F_{u_1u_2}$ and
$\Gamma_{13}\subset F_{u_1u_3}$; by (T3) the corner of $S$ lies in $F_{u_1}$. Since $\{0,u_1,u_2,u_3\}$ is a unimodular simplex,
there are $\nu_2,\nu_3\in N$ with $\langle u_1,\nu_k\rangle=0$ and $\langle u_j,\nu_k\rangle=\delta_{jk}$ for $j,k\in\{2,3\}$.
Choose $z_1,z_2$ as above with $z_1z_2=x^{-u_1}$ and $\langle m_i,\nu_2\rangle=\langle m_i,\nu_3\rangle=0$, and put
$\mathfrak t_j=x^{u_j-u_1}$. In $U_\sigma$, $W$ is $\beta_0z_1z_2=\alpha_{u_1}+\alpha_{u_2}\mathfrak t_2+\alpha_{u_3}\mathfrak
t_3$, and $\mathrm{Crit}=\{z=0,\ \alpha_{u_1}+\alpha_{u_2}\mathfrak t_2+\alpha_{u_3}\mathfrak t_3=0\}$ is a localized line in the
$(\mathfrak t_2,\mathfrak t_3)$-torus, a pair of pants whose three ends lie over the three legs of $\widehat\Gamma$ at $p$. Along
$\Gamma_{12}$ the leg collar paths of the two-face $F_{u_1u_2}$ arrive, by \eqref{eq:collarphase}, with lifted phases
\[
(\arg\mathfrak t_2,\arg\mathfrak t_3)=\bigl(\epsilon_{12}\pi,\ \langle u_3-u_1,\nu_{12}\rangle\pi\bigr),\qquad
\epsilon_{12}=\langle u_2-u_1,\nu_{12}\rangle=\pm1,
\]
where $\nu_{12}$ is the cocharacter of $F_{u_1u_2}$, and $x^{-u_1}$ ends in $\R_{\ge0}$, since the side is that of $F_{u_1}$;
symmetrically along $\Gamma_{13}$. In the part of a neighbourhood of $\Gamma_{12}$ where $u_3$ is inactive, the circle
$\nu_3\otimes U(1)$ preserves $W$ (Lemma~\ref{lem:symm}), commutes with $S^1_v$, preserves $\mathrm{Crit}$, and changes, among
$\mathfrak t_2,\mathfrak t_3,x^{-u_1}$, only the phase of $\mathfrak t_3$. Let $\mathbb B_p$ be the corner of $S$ at $p$, a closed
disc bounded by segments $[p,e_2]\subset\Gamma_{12}$ and $[p,e_3]\subset\Gamma_{13}$ of the two legs, the transverse segments at
$e_2$ and $e_3$, and an arc of the inner boundary; the segments are long compared with $\varsigma_0$, so that the collar paths over
points of $\Gamma_{12}$ near $e_2$ lie in the part where $u_3$ is inactive, and symmetrically at $e_3$. The \emph{slide} along
$\Gamma_{12}$ modifies the collar paths over a segment $[e'_2,e_2]\subset\Gamma_{12}$ in that part: it replaces
$\gamma_e(\varsigma)$ by $e^{i\pi k\phi(e)\chi(\varsigma)\nu_3}\gamma_e(\varsigma)$, where $\phi$ increases from $0$ at $e'_2$
to $1$ at $e_2$, $\chi=1$ for $\varsigma\le\varsigma_0/4$ and $\chi=0$ for $\varsigma\ge\varsigma_0/2$, and the integer $k$ is
chosen so that the lifted phase of $\mathfrak t_3$ at the end becomes $\epsilon_{13}\pi$, the value with which the collar paths
along $\Gamma_{13}$ arrive. Symmetrically, the slide along $\Gamma_{13}$ turns $\arg\mathfrak t_2$ to $\epsilon_{12}\pi$ with
$\nu_2$. After the slides both families arrive at $e_2$ and $e_3$ with $(\arg\mathfrak t_2,\arg\mathfrak
t_3)=(\epsilon_{12}\pi,\epsilon_{13}\pi)$ and $x^{-u_1}\in\R_{\ge0}$, that is, in the \emph{corner chamber}
$W\cap U_\sigma\cap\{x^{-u_1}\in\R_{\ge0},\ \mathfrak t_2,\mathfrak t_3\in\R_{<0}\}$, which contains the real arc of $\mathrm{Crit}$
joining $\Gamma_{12}$ to $\Gamma_{13}$; on it $x^{-u}$, for the first element $u$ of the ordering of $F_{u_1u_j}$, lies in
$s\R_{\ge0}$, $s$ being the side of the corner sheet at $\Gamma_{1j}$. The piece is $S^1_v\cdot\Theta_p$ over
$\mathbb B_p$, with $\Theta_p$ given by Lemma~\ref{lem:collar}(3), together with the slid leg pieces; near $p$, $\Theta_p$ may be
taken with values in the corner chamber, as in the model of \cite[\S7.4]{CBM}, where the swept set is homeomorphic to
$\C\times\R$.

\emph{Crossings.} At a crossing $\mathfrak x_k$ of a leg with covector $\alpha=u'-u$, $v$ is monodromy invariant, that is
$v\in(u,u')^\perp$, so $S^1_v$ acts exactly on the local model (Lemma~\ref{lem:symm}). The piece is $S^1_v\cdot\mathfrak s(D_k)$ for a
small disc $D_k$ around $\hat j(\mathfrak x_k)$ in $\hat S$, a solid torus; $\mathfrak s$ lies in the dense torus, where $S^1_v$ acts
freely.

\emph{Nodes.} Let $q$ be a node on $\partial S$, with parallelogram $P=abcd$ and $F_q=F_{u_+u_-}$. Near
$C^W_q:=W\cap V(\operatorname{cone}(a,c))$ only the monomials $0,u_\pm$ can be active, because the dual edge $[u_-,u_+]$ has
length one, and the constant term is cut off near the strata; hence near $C^W_q$, $W$ is
$\{\mathfrak t=\mathfrak t_0\}\times\widehat{\mathrm{Con}}_q$, where $\mathfrak t=x^{u'-u}$ for the ordering $(u,u')$ of $F_q$ and
$\widehat{\mathrm{Con}}_q$ is the resolved conifold of the two cones $\operatorname{cone}(a,b,c)$, $\operatorname{cone}(a,c,d)$ in
$N_q=\operatorname{span}(a,b,c,d)\cap N$; $N=N_q\oplus\Z\nu$ for the cocharacter $\nu$ of $F_q$, and $N_q\otimes U(1)$ acts
exactly. So $C^W_q\cong\mathbb P^1$ is the exceptional curve of $W$ at $q$. The lattice
$\Lambda_f=(u_+,u_-)^\perp\cap N=\Z(b-a)\oplus\Z(c-b)$ is the lattice of differences in $N_q$. On $C^W_q$, with coordinate
$\zeta=Z_b/Z_d$, the circle $S^1_\lambda$, $\lambda=k_1(b-a)+k_2(c-b)$, acts by $\zeta\mapsto e^{i(k_1-k_2)\theta}\zeta$; we call
$\mathrm{wt}(\lambda)=k_1-k_2$ its \emph{weight}. The curves $D_{ab}=\{Z_a=Z_b=0\}$ and $D_{bc}$ meet $C^W_q$ at $\zeta=0$, and
$D_{cd}$, $D_{da}$ at $\zeta=\infty$. Let $\mathbb B_q\subset\mathfrak H_q^\circ$ be a closed disc bounded by the part of the
boundary path of $\mathfrak H_q$ between points $e'$, $e''$ of the two leg segments near $y_0$ and $y_\infty$, the transverse
segments of $\mathfrak H_q^\circ$ at $e'$ and $e''$, and an arc of the inner boundary. The \emph{node piece} is $S^1_v\cdot\Theta_q$ over
$\mathbb B_q$, together with the leg pieces over the rest of the two leg segments, where $\Theta_q\colon\mathbb B_q\to W$ is given
by Lemma~\ref{lem:collar}(1): it agrees with the leg collar paths $\gamma_{e'}$, $\gamma_{e''}$ on the transverse segments and
with $\mathfrak s$ on the inner boundary, maps the boundary path near $y_0$ and $y_\infty$ to $\mathrm{Crit}$, and maps $e_q$ to an arc
$\gamma_q\subset C^W_q$ from $\zeta=0$ to $\zeta=\infty$.

\begin{lemma}[node degree]\label{lem:nodedeg}
Let $q$ be a node on $\partial S$. The node piece restricts on $\partial_q\mathfrak D$ to a map onto $C^W_q$ of degree
$-\rho_q(S)$, for the complex orientation of $C^W_q$. More generally, if $v_q=-(c_{ab}(b-a)+c_{bc}(c-b))$ is the coefficient on
$e_q$ (oriented from $y_0$ to $y_\infty$) determined by balancing at $y_0$, where $c_{ab}(b-a)$ and $c_{bc}(c-b)$ are the
coefficients of the legs with lower supports $\{a,b\}$ and $\{b,c\}$, oriented away from the node as in
Section~\ref{ssec:cycles-bdry}, the degree is $\mathrm{wt}(v_q)=c_{bc}-c_{ab}$.
\end{lemma}
\begin{proof}
On $\partial_q\mathfrak D$ the map is $(x,\theta)\mapsto e^{i\theta v}\cdot\gamma_q(x)$, $x\in e_q$, collapsed at the poles,
which $S^1_v$ fixes. Take the $b-a$ passage entering along the leg with lower support $\{a,b\}$, with $v=\epsilon_v(b-a)$. The
boundary path of $\mathfrak H_q$ runs from that leg through $y_0$, along $e_q$ and through $y_\infty$ to the leg with lower
support $\{c,d\}$, and $\gamma_q$ runs from $\zeta=0$ (on $D_{ab}$) to $\zeta=\infty$ (on $D_{cd}$). By Stokes, the orientation of
$\partial_q\mathfrak D$ is (boundary orientation of $e_q$)$\times$(circle), and in polar coordinates $\zeta=re^{i\phi}$ the map is
$(r,\theta)\mapsto re^{i\epsilon_v\theta}$, of degree $\epsilon_v=\mathrm{wt}(v)$. Since $c_{ab}=-\epsilon_v$ and $c_{bc}=0$, this is
$\mathrm{wt}(v_q)$, and it equals $-\rho_q$ by the formula $\rho_q=c_{ab}-c_{bc}$ of Section~\ref{ssec:cycles-bdry}. For the
$c-b$ passage entering along the leg with lower support $\{b,c\}$, with $v=\epsilon_v'(c-b)$, the arc again runs from $\zeta=0$ to
$\zeta=\infty$ and the degree is $\mathrm{wt}(v)=-\epsilon_v'=-\rho_q$. Reversing the orientation of $\partial S$ reverses both
$\gamma_q$ and the sign of $\rho_q$. The degree depends only on the homotopy class of $\gamma_q$ relative to the poles, and
$C^W_q$ is simply connected, so it does not depend on the choice of $\gamma_q$.
\end{proof}

The existence of the maps $\Theta$ rests on a computation of local fundamental groups, proved in Appendix~\ref{app:collar}.
Let $(\tau,\omega)$ be a dual pair, $u_0$ a vertex of $\tau$ and $\tau'=\tau\setminus\{u_0\}$, and let
$X_\omega\subset X_{\mathcal F'}$ be the open toric subvariety of the cones of $\mathcal F'$ contained in
$\operatorname{cone}(\omega)$. The \emph{extended tropicalisation} of $X_\omega$ is $\operatorname{Log}$ on the dense torus and,
on the orbit $O(\sigma)$ of a cone $\sigma\subseteq\operatorname{cone}(\omega)$, the logarithm of that orbit, with values in
$N_\R/\operatorname{span}(\sigma)$; a subset $\Omega\subseteq N_\R$ has a preimage under it, consisting of the points whose
logarithm lies in $\Omega$, respectively in the image of $\Omega$ in $N_\R/\operatorname{span}(\sigma)$. The four cases used
are: a node ($\tau=[u_+,u_-]$, $\omega=P$, subdivided by its diagonal in $\mathcal F'$), a negative vertex ($\tau$ an edge,
$\omega$ a triangle), a positive vertex ($\tau$ a triangle, $\omega$ an edge) and a switching point ($\tau$ and $\omega$ edges).

\begin{lemma}[local fundamental groups]\label{lem:pi1}
Let $\dim\tau+\dim\omega=3$, or $\dim\tau=\dim\omega=1$, and let $N_\omega=\operatorname{span}(\omega)\cap N$. The characters
$u-u_0$, $u\in\tau'$, are part of a basis of the saturated lattice $N_\omega^\perp\subset M$; if $\dim\tau+\dim\omega=3$ they
are a basis, and if $\dim\tau=\dim\omega=1$ we complete them to a basis by a character $\kappa$. Let $\nu_u\in N$, $u\in\tau'$,
and in the second case $\nu_\kappa\in N$, satisfy $\langle u_0,\cdot\rangle=0$ and pair with this basis of $N_\omega^\perp$ to
the dual basis; so $\langle u',\nu_u\rangle=\delta_{uu'}$ for $u,u'\in\tau'$. Then $N=N_\omega\oplus\bigoplus_{u\in\tau'}\Z\nu_u$,
respectively $N=N_\omega\oplus\Z\nu_u\oplus\Z\nu_\kappa$. Let $\Omega=\Omega'\oplus\sum_{u\in\tau'}I_u\nu_u$, respectively
$\Omega=\Omega'\oplus I_u\nu_u\oplus I_\kappa\nu_\kappa$, with $\Omega'\subset N_{\omega,\R}$ open and convex,
$\Omega'+\operatorname{cone}(\omega)=\Omega'$, and $I_u$, $I_\kappa$ bounded open intervals. Assume that only the monomials $0$
and $u\in\tau$ are active on $\operatorname{Log}^{-1}(\Omega)$, and that $\Omega$ contains the ball of radius $4C_0$ about a point
of $\operatorname{ri}F_\tau$. Let $W_\Omega$ be the preimage of $\Omega$ in $W\cap X_\omega$ under the extended tropicalisation.
Then $W_\Omega$ is open in $W$ and path connected, and the characters $\mathfrak t_u=x^{u-u_0}$, $u\in\tau'$, induce an
isomorphism $\pi_1(W_\Omega)\to\Z^{\tau'}$, respectively, together with $w=x^\kappa$, an isomorphism
$\pi_1(W_\Omega)\to\Z^{\tau'}\times\Z$.
\end{lemma}

\begin{lemma}[consistency of the collar paths]\label{lem:collar}
\begin{enumerate}
\item[(1)] At a node $q$ on $\partial S$, for every arc $\gamma_q\subset C^W_q$ from $\zeta=0$ to $\zeta=\infty$, the map on
$\partial\mathbb B_q$ given by the leg collar paths $\gamma_{e'}$, $\gamma_{e''}$, by $\mathfrak s$ on the inner boundary arc, by
$\mathrm{Crit}$ over the leg segments and by $\gamma_q$ over $e_q$ extends to a continuous map $\Theta_q\colon\mathbb B_q\to W$.
\item[(2)] At a negative vertex, on each of the three sheet corners $\mathbb B_j$, the map on $\partial\mathbb B_j$ given by the
leg collar paths, $\mathfrak s$, $\mathrm{Crit}$ and the vertex path extends to a continuous map $\Theta_j\colon\mathbb B_j\to W$.
\item[(3)] At a positive vertex, after the slides, the map on the boundary of the corner $\mathbb B_p$ given by the leg collar
paths, $\mathfrak s$ and the real arc of $\mathrm{Crit}$ extends to a continuous map $\Theta_p\colon\mathbb B_p\to W$.
\item[(4)] At a switching point $e^*$, the map on $\partial\mathbb B_{e^*}$ given by the leg collar paths $\gamma_{e'}$,
$\gamma_{e''}$, by $\mathfrak s$ on the inner boundary arc and by $\mathrm{Crit}$ over $[e',e'']$ extends to a continuous map
$\Theta_{e^*}\colon\mathbb B_{e^*}\to W$.
\end{enumerate}
Explicitly, in (1) the two families of collar paths end, in the conifold coordinates $z_1=Z_aZ_b$, $z_2=Z_cZ_d$, $z_3=Z_aZ_d$,
$z_4=Z_bZ_c$ taken as characters vanishing on $\nu$, with $(\operatorname{sign}z_1,\operatorname{sign}z_2)=(s_q,+)$ on
the leg with lower support $\{a,b\}$ and $(+,s_q)$ on the leg with lower support $\{c,d\}$ for the $b-a$ passage, and
with $(\operatorname{sign}z_4,\operatorname{sign}z_3)=(s_q,+)$ on $\{b,c\}$ and $(+,s_q)$ on $\{d,a\}$ for the
$c-b$ passage. At side $s_q=+1$ the ends lie in one chamber, and near $C^W_q$ the map $\Theta_q$ can be taken with
values in the positive real part of the local model $\{\mathfrak t=\mathfrak t_0\}\times\widehat{\mathrm{Con}}_q$, the union of
whose $S^1_v$-orbits is the set $\pi^{-1}(P_0)$ of \cite[Lemma~4.2]{CBM}, with boundary the exceptional curve; $\gamma_q$ is then
the positive real arc of $C^W_q$. At side $s_q=-1$ the element $e^{i\pi\lambda}$, with $\lambda=c-b$ for the $b-a$
passage and $\lambda=a-b$ for the $c-b$ passage, acts by $(-1,-1)$ on $(z_1,z_2)$, respectively $(z_4,z_3)$, and exchanges the
two chambers; $\Theta_q$ can be taken with $\gamma_q$ the positive real arc followed by the half-turn $e^{i\pi\phi(x)\lambda}$
along $e_q$, $\phi$ increasing from $0$ at $y_0$ to $1$ at $y_\infty$.

None of these choices (the cocharacters $\nu$, the slides, the switching discs, $\gamma_q$ and the extensions) affects the
boundary computation: $\partial F_*[\mathfrak D]$ depends on them only through the maps $F|_{\partial_q\mathfrak D}$, whose
degrees are given by Lemma~\ref{lem:nodedeg}.
\end{lemma}
\begin{proof}
In each case the domain ($\mathbb B_q$, $\mathbb B_j$, $\mathbb B_p$, $\mathbb B_{e^*}$) is a disc, and the boundary data form a
loop in the set $W_\Omega$ of Lemma~\ref{lem:pi1} for a suitable $\Omega$. Let the dual pair $(\tau,\omega)$ be
\[
([u,u'],P),\qquad([u,u'],\{n_1,n_2,n_3\}),\qquad(\{u_1,u_2,u_3\},[n,n']),\qquad([u,u'],[n,n'])
\]
at a node, a negative vertex, a positive vertex and a switching point respectively, with $u_0=u$, respectively $u_1$. Then the $\nu_u$ may be taken to be the cocharacters $\nu$, respectively
$\nu_2,\nu_3$, of the adapted coordinates, and at a switching point $\kappa$ is the character of the coordinate $w$ and
$\nu_\kappa$ the fourth element of the dual basis of $(m_1,m_2,u'-u,\kappa)$.

\emph{Sizes, step 1: the sets.} Let $\mathcal X_0\subset\operatorname{ri}F_\tau$ be the compact set $y_0\cup e_q\cup y_\infty$,
$\{y\}$, $\{p\}$ or $\{e^*\}$, and let $\mathcal X$ be the union of $\mathcal X_0$ with the images under $\hat j$ of the arcs of
$\partial S$ and of the interior edge in the boundary of the domain. Since $\operatorname{conv}\mathcal X_0$ is a compact subset
of $\operatorname{ri}F_\tau$, there is $d_0>0$ such that on its $d_0$-neighbourhood every monomial other than $0$ and the vertices
of $\tau$ lies at least $d_0$ below the maximum of the $T_i$. The lengths of these arcs (the corner lengths and the widths of the
switching discs), at most $\varsigma_1$, and the transverse length $\varsigma_0$ are chosen with $\mathcal N_\Gamma$, with
$\varsigma_0$ small compared with $\varsigma_1$ and $\varsigma_1$ small compared with $d_0$.

\emph{Sizes, step 2: where the boundary data lie.} The section $\mathfrak s$ over the inner boundary arc and the transverse
segments, and the collar paths, lie over the $d_1(\varsigma_0+C_0t_{\max})$-neighbourhood of $\mathcal X$ plus
$\operatorname{cone}(\omega)$, with $d_1$ depending only on the adapted coordinates: $\mathfrak s$ lies in the band of
Lemma~\ref{lem:pos}, and steps (ii)--(iv) change only $|\mathfrak t|$ and $|z_1z_2|$, which moves the logarithm by
$O(\varsigma_0+C_0t_{\max})$ in the direction of $\nu$ and otherwise along $\operatorname{cone}(\omega)$.

\emph{Sizes, step 3: the set $\Omega$.} Let $\widetilde{\mathcal X}$ be the open
$\bigl(d_1(\varsigma_0+C_0t_{\max})+4C_0\bigr)$-neighbourhood of $\operatorname{conv}\mathcal X$, an open convex set within
$\varsigma_2=\varsigma_1+d_1(\varsigma_0+C_0t_{\max})+4C_0$ of $\operatorname{conv}\mathcal X_0$. Let $\Omega'$ be the projection
of $\widetilde{\mathcal X}$ to $N_{\omega,\R}$ along the $\nu_u$ (and $\nu_\kappa$), plus $\operatorname{cone}(\omega)$, and $I_u$
and $I_\kappa$ the projections of $\widetilde{\mathcal X}$ to the lines of $\nu_u$ and $\nu_\kappa$. The $\nu_u$-component of a
point $y$ is $\langle u-u_0,y\rangle$, which is constant on $\operatorname{aff}F_\tau\supset\operatorname{conv}\mathcal X_0$, and
at a switching point $\mathcal X_0$ is a point; so $\Omega$ lies within $d_2\varsigma_2$ of
$\operatorname{conv}\mathcal X_0+\operatorname{cone}(\omega)$, with $d_2$ depending only on the splitting.

\emph{Sizes, step 4: the active monomials.} For $\varsigma_1$ small compared with $d_0$ and $C_0$ small, only $0$ and the
vertices of $\tau$ are active on the $d_2\varsigma_2$-neighbourhood of $\operatorname{conv}\mathcal X_0$, and this persists on
its sum with $\operatorname{cone}(\omega)$, because for $z\in\operatorname{cone}(\omega)$, $u\in\tau$ and $m\in\Delta^\circ$ the
difference $T_u-T_m$ changes by $\langle m-u,z\rangle\ge0$ from $y$ to $y+z$.

\emph{Sizes, step 5: conclusion.} So $\Omega$ satisfies the hypotheses of Lemma~\ref{lem:pi1}, and the boundary data lie in
$W_\Omega$: $\mathfrak s$ and the collar paths by the above, and $\mathrm{Crit}$ and $C^W_q$ because they lie on the strata of
$X_\omega$ over $\mathcal X$. By Lemma~\ref{lem:pi1} the loop is null-homotopic in $W_\Omega$, and the map extends over the disc,
if and only if the total change of the lifted phases of the $\mathfrak t_u$, and at a switching point of $w$, around the loop is
zero.

The switching points are finitely many and lie away from the vertices and the nodes (Section~\ref{ssec:res-shell}), so we choose
the corner lengths, the slide segments and the switching discs so that no corner arc, no slide segment and no arc of
$\partial S$ in the boundaries of $\mathbb B_q$, $\mathbb B_j$ and $\mathbb B_p$ contains a switching point, and so that the
switching discs are disjoint from the slide segments and from these discs. In particular the collar paths are unslid over every
switching disc, so they keep $w$ fixed there, as (4) uses. Moreover the collar paths on the transverse segments bounding these
discs then have the side of the corresponding sheet at the vertex or node; at a node this holds because
$\mathbb B_q\subset\widehat U_q$.

(1) Here $\tau'=\{u'\}$ and $\mathfrak t=x^{u'-u}$. Around $\partial\mathbb B_q$: along $\mathfrak s$ the phase of $\mathfrak t$
is $0$; each leg collar path in $F_q$ uses the cocharacter $\nu$ of $F_q$ and increases the lifted phase of $\mathfrak t$ by
$\pi$ from $\mathfrak s$ to $\mathrm{Crit}$, by \eqref{eq:collarphase}; and on $\mathrm{Crit}$ and on $C^W_q$ we have
$\mathfrak t=\mathfrak t_0<0$, so the phase is constantly $\pi$. Going around, the change is $0+\pi+0-\pi=0$, whatever the sides
and the arc $\gamma_q$. The chamber pattern follows from \eqref{eq:collarphase}: on the leg with lower support $\{a,b\}$ the
adapted coordinates are $z_1^{\rm leg}=z_3$, $z_2^{\rm leg}=z_4$, $w=z_2$, so the collar path ends with $z_2,z_4>0$ and
$z_3,\ z_1=z_3z_4/z_2\in s_q\R$; on the leg with lower support $\{c,d\}$ the roles of $z_1,z_2$ and of $z_3,z_4$ are
exchanged; the $c-b$ passage is the same with $w=z_3$ on $\{b,c\}$ and $w=z_4$ on $\{d,a\}$. The element $e^{i\pi(c-b)}$ has
weights $\langle(1,1,0,0),c-b\rangle=-1$ and $\langle(0,0,1,1),c-b\rangle=1$ on $z_1,z_2$ and $0$ on $z_3,z_4$ (vectors of values
on $a,b,c,d$), and similarly for $e^{i\pi(a-b)}$; both lie in $\Lambda_f\otimes U(1)$, which acts exactly near $C^W_q$ and rotates
$C^W_q$ with weight $\mp1$.

(2) Here too $\tau'=\{u'\}$. Around $\partial\mathbb B_j$ the lifted phase of $\mathfrak t$ changes by $0$ along $\mathfrak s$, by
$+\pi$ along the leg collar path, by $0$ on $\mathrm{Crit}$, where it is constantly $\pi$, and by $-\pi$ back along the vertex
path, which also turns $\arg\mathfrak t$ by $\pi$. On the interior edge $\Theta_j$ is the vertex path for each $j$, so
$G_j\cdot\Theta_j$ and the cone on $\mathcal Q$ agree there, $T^2$ acting exactly and $\partial\mathcal Q=G_1+G_2+G_3$.

(3) Here $\tau'=\{u_2,u_3\}$. Around $\partial\mathbb B_p$ the lifted phases of $(\mathfrak t_2,\mathfrak t_3)$ change by $(0,0)$
along $\mathfrak s$, by $(\epsilon_{12},\epsilon_{13})\pi$ along the slid collar path at $e_2$ from $\mathfrak s$ to $\mathrm{Crit}$, by
$(0,0)$ along the real arc of $\mathrm{Crit}$, on which $\mathfrak t_2,\mathfrak t_3<0$, and by $-(\epsilon_{12},\epsilon_{13})\pi$ back
along the slid collar path at $e_3$. The total is zero.

(4) Here $\tau'=\{u'\}$, and $\pi_1(W_\Omega)\cong\Z^2$ is detected by $(\mathfrak t,w)$. Around $\partial\mathbb B_{e^*}$ the
lifted phase of $\mathfrak t$ changes by $0$ along $\mathfrak s$, by $+\pi$ along $\gamma_{e'}$ from $\mathfrak s$ to
$\mathrm{Crit}$, whatever the side at $e'$, by $0$ along $\mathrm{Crit}$, where $\mathfrak t=\mathfrak t_0$, and by $-\pi$ back
along $\gamma_{e''}$; the total is $0$. The lifted phase of $w$ does not change: $w>0$ on $\mathfrak s$ and on $\mathrm{Crit}$
over $[e',e'']$, and the collar paths keep $w$ fixed.

For the last statement: by Lemma~\ref{lem:domain} the boundary of $[\mathfrak D]$ is supported on the spheres
$\partial_q\mathfrak D$, on which $F$ is $(x,\theta)\mapsto e^{i\theta v}\gamma_q(x)$, whose degree does not depend on $\gamma_q$
(Lemma~\ref{lem:nodedeg}).
\end{proof}

Without the slides, the lifted phases around $\partial\mathbb B_p$ would change by
$\bigl(\epsilon_{12}-\langle u_2-u_1,\nu_{13}\rangle,\ \langle u_3-u_1,\nu_{12}\rangle-\epsilon_{13}\bigr)\pi$, which is in general
nonzero; the slides are therefore needed, and their amounts are fixed by the cocharacters. At nodes and negative vertices no
correction of the phase of $\mathfrak t$ is needed, because all collar paths there lie in one two-face and use one cocharacter;
the half-turn at a node on side $-1$ is part of an explicit choice of the arc $\gamma_q$, not a condition for the existence
of $\Theta_q$, which holds for every arc. The
normalisation $\langle u,\nu\rangle=0$ is what allows adapted coordinates with $z_1z_2=x^{-u}$ and
$\langle m_1,\nu\rangle=\langle m_2,\nu\rangle=0$; it is used for the explicit chamber pattern in Lemma~\ref{lem:collar}, while
the existence of the maps $\Theta$ uses only that all collar paths in one two-face turn $\arg\mathfrak t$ by the same $+\pi$.

The pieces satisfy the following \emph{inner boundary condition}:
\begin{enumerate}
\item[(I)] on the boundary of $\mathcal N_\Gamma$ in $S$, the local pieces are $S^1_v$-orbits through points $\mathfrak s(b)$, and on
the two-cells $\mathcal Q\cdot\mathfrak s(b)$ at the interior edges leaving boundary vertices they lie in
$T^2\cdot\mathfrak s(b)$; in both cases the base points $b$ lie in the set $V_c$ of the adjacent cell $c$ of $\mathfrak T$, so the
images lie in $W_{V_c}$.
\end{enumerate}

\subsection{The regular part}\label{ssec:res-ext}
Over $S\setminus\operatorname{int}\mathcal N_\Gamma$ the map is extended cell by cell, each cell being mapped into the region
of an open set of $\widehat B$ that is product-shaped, contains the image of the cell, and increases with the cell.

\begin{lemma}[admissible cover]\label{lem:AC}
For the choices listed in Appendix~\ref{app:AC}, there are a triangulation $\mathfrak T$ of $S\setminus\operatorname{int}\mathcal
N_\Gamma$, adapted to the faces of $\nabla^{\check h}$ (each closed cell meets each closed face in one face of the cell), and for
each cell $c$ an open set $V_c\subset\widehat B$ such that:
\begin{enumerate}
\item $\hat j(c)\subset V_c$, and $V_c$ is product-shaped at the face $\Phi_c$ of lowest dimension met by $\hat j(c)$, in a chart
of $(\Phi_c,n_c)$ for a vertex $n_c$ whose vertex region contains $\hat j(c)\cap\Phi_c$;
\item if $c'\subset c$ then $\Omega_{V_{c'}}\subset\Omega_{V_c}$, hence $W_{V_{c'}}\subset W_{V_c}$;
\item for every $r>0$ the choices can be made so that $V_c$ lies in the $r$-neighbourhood of $\hat j(c)$.
\end{enumerate}
\end{lemma}
The proof is in Appendix~\ref{app:AC}. Its main point is that near an edge crossing $x_0$ the cells must be graded (the
transverse radius of $V_c$ depends on the lowest face of $c$) and the ring of cells around $x_0$ must be a cone from $x_0$;
admissibility is reduced to a strict inequality on a compact set, which holds by conicity of $\hat S$ and of $\widehat B$ at
$x_0$.

\begin{lemma}[extension]\label{lem:ext}
Given local pieces satisfying (I), the map on $\partial\mathfrak D_{\rm reg}$ extends to a map $F$ on
$\mathfrak D_{\rm reg}=\mathfrak D|_{S\setminus\operatorname{int}\mathcal N_\Gamma}$ that sends the cells over each cell $c$ of
$\mathfrak T$ into $W_{V_c}$.
\end{lemma}
\begin{proof}
Build $F$ by increasing dimension of cells, mapping the cells over each cell $c$ of $\mathfrak T$ into the region $W_{V_c}$. By
Lemma~\ref{lem:AC}(2), the cells over the faces of $c$ are then already mapped into $W_{V_c}$. On $0$-cells use $\mathfrak s$. On
$1$-cells in the base direction use $\mathfrak s$; on circle fibres use $S^1_v$-orbit arcs through $\mathfrak s$ where $S^1_v$ acts
exactly, and otherwise arcs whose loops represent $v\in\Lambda_{V_c}=\pi_1(W_{V_c})$ (Lemma~\ref{lem:regions}). A map on the
boundary of a $2$-cell extends if and only if the boundary loop is trivial in $\pi_1(W_{V_c})$, which is abelian, so base points
do not matter. A section$\times$section cell is filled by $\mathfrak s(c)$, and its boundary is contractible by
Lemma~\ref{lem:pos}. For a section$\times$fibre cell the two fibre loops represent the same element $v$ of $\Lambda_{V_c}$, by
Lemma~\ref{lem:regions}(2) and because $v$ is parallel. For the triangles $\mathcal Q$ over interior edges the boundary loop is
$s_1v_1+s_2v_2+s_3v_3=0$, which is clause (g). At an interior vertex the four triangles are flat, so their union is the boundary
of the image of the flat $3$-simplex, and the classes agree across the four edges. For cells of dimension at least three the
extension exists because $W_{V_c}$ is aspherical.
\end{proof}

\subsection{Transport}\label{ssec:res-transport}
An orbit-preserving isotopy carries the constructed chain from the localized model to a holomorphic member.
Transport through the nondegenerate family then identifies this member with the chosen small resolution, preserving the exceptional curves.

\begin{lemma}\label{lem:T0}
For $R$ large there is a homeomorphism $\Upsilon^0_{\rm res}$ of $X_{\mathcal F'}$, isotopic to the identity, smooth on the smooth locus
and preserving every torus orbit, with $\Upsilon^0_{\rm res}(W)=\widehat X_R$ and, for every $q$,
$\Upsilon^0_{\rm res}(C^W_q)=C_q^R$ with degree one, where $C_q^R:=\widehat X_R\cap V(\operatorname{cone}(a,c))$.
\end{lemma}
\begin{proof}
Let $\partial_1,\dots,\partial_8$ be the fundamental vector fields of the complex torus $N\otimes\C^*$, regarded as a real Lie
group: for a basis $e_1,\dots,e_4$ of $N$, the radial fields generating $r\mapsto e^{r}\cdot x$ along $e_j\otimes\R_{>0}$ and the
angular fields generating $\theta\mapsto e^{i\theta e_j}\cdot x$. They are smooth on $X_{\mathcal F'}$, since the action is
algebraic, tangent to every orbit, and they span the tangent space of each orbit, on which the torus acts transitively. Let
$\mathcal G(x,s)$ be the defining function of $V_s$ in local Cox coordinates. By Lemma~\ref{lem:Ufan}(3), at every point $(x,s)$
of $\mathcal V=\bigcup_sV_s\times\{s\}$ the real derivatives $d\mathcal G(\partial_j)$ span $\C$; by continuity two of them,
$d\mathcal G(\partial_i)$ and $d\mathcal G(\partial_k)$, are $\R$-linearly independent on a neighbourhood, and there
$\mathfrak X=\partial_s+a_i\partial_i+a_k\partial_k$, with $(a_i,a_k)$ the smooth solution of
$a_i\,d\mathcal G(\partial_i)+a_k\,d\mathcal G(\partial_k)=-\partial_s\mathcal G$, is tangent to $\mathcal V$. A partition of
unity glues these local fields, since the condition is affine in the coefficients, and the coefficients are cut off outside a
compact neighbourhood of $\mathcal V$ in the smooth locus (Lemma~\ref{lem:Ufan}(2)). The field $\mathfrak X$ is tangent to
$O(\sigma)\times[0,1]$ for every $\sigma$, so its flow preserves orbits; its $\partial_s$-component is $1$, so the flow from $s=1$
to $s=0$ exists, and it carries $V_1=W$ to $V_0=\widehat X_R$. (The radial fields are needed: at a positive real point in a
cut-off band, $\partial_s\mathcal G$ is real and nonzero while the derivatives along the angular fields are purely imaginary.) On
$V(\operatorname{cone}(a,c))$ only $u_\pm$ survive, and the restricted defining function depends only on the character
$z=x^{u_+-u_-}$, the cut-offs only on $|z|$; so $V_s\cap V(\operatorname{cone}(a,c))$ is the closure of a level set $\{z=z_s\}$
of the orbit, $z_s$ the unique regular zero of the radial equation, and these curves form an isotopy. The flow preserves
$V(\operatorname{cone}(a,c))$, hence maps $C^W_q$ onto $C^R_q$; its degree is $+1$ because $\Upsilon^0_{\rm res}$ is isotopic to the
identity and $D_b\cdot C_q=1$.
\end{proof}

\begin{lemma}\label{lem:T}
Let $\mathcal R$ be the set of coefficient vectors of $\Delta^\circ$-regular hypersurfaces, a connected Zariski open set. For
every $X$ in $\mathcal R$ and every path in $\mathcal R$ from $X_R$ to $X$ there is a homeomorphism $\Upsilon^1_{\rm res}$ of
$X_{\mathcal F'}$, isotopic to the identity, smooth on the smooth locus and preserving every torus orbit, which carries
$\widehat X_R$ to $\widehat X$ and $C_q^R$ to $C_q$ with degree one.
\end{lemma}
\begin{proof}
$X_R\in\mathcal R$ by Lemma~\ref{lem:T0} at $s=0$. For $f\in\mathcal R$ the proper transform $\widehat X_f$ is transverse to
every orbit of $X_{\mathcal F'}$, because each orbit maps to an orbit of $\mathbb P_\Delta$ by a surjective homomorphism of tori,
which is a submersion. The construction of Lemma~\ref{lem:T0} along the path gives $\Upsilon^1_{\rm res}$. It preserves
$V(\operatorname{cone}(a,c))$, hence carries $C^R_q$ to $C_q$. By Lemma~\ref{lem:Ufan}(1) and the choice of $\mathcal F'$, the
proper transform of $X$ in $X_{\mathcal F'}$ is the resolution $\widehat X$.
\end{proof}

Put $\Upsilon_{\rm res}=\Upsilon^1_{\rm res}\circ\Upsilon^0_{\rm res}$. The path is one of the choices; the relative class of $\Xi(S)$ may change with it by an
element of $H_3(\widehat X)$.

\subsection{Proof of Theorem~\ref{thm:res}}\label{ssec:res-proof}
We assemble the local maps and the extension over the regular part, with the orientations fixed above.
The degree computation at each node gives the boundary coefficient, and transport yields the chain on the small resolution.

Let $(S,j,v)$ be a tropical $2$-cycle on $B'_\sigma$, with $j$ piecewise linear, and let $\widehat X$ and $\mathcal F'$ be as in
Section~\ref{ssec:res-base}, the node coefficients being computed with the diagonals of $\widehat X$. The choices are made in the
following order, each depending only on the previous ones:
\begin{enumerate}
\item the orderings and cocharacters $\nu$ of the two-faces, the characters $\kappa$ and the adapted coordinates at the legs,
and the cocharacters used for the slides (these are combinatorial);
\item the standard position $j'$, the sides $s_q$ and $r_*$ (Lemma~\ref{lem:NS});
\item the homeomorphism $g$ and the half-discs $\mathfrak H_q$ (Lemma~\ref{lem:resmove}, Definition~\ref{def:resolvedcycle});
\item the normalising isotopy (T1)--(T3), which determines the switching points; then the neighbourhood $\mathcal N_\Gamma$,
with the switching discs, the corner lengths and the transverse length $\varsigma_0$ as in the proof of Lemma~\ref{lem:collar};
and the generic isotopy (Section~\ref{ssec:res-shell});
\item the charts, the triangulation $\mathfrak T$, the radii and $\eta$ (Appendix~\ref{app:AC});
\item $C_0$ small, the cut-off profile $\beta$ and the constant $C_1$, and then $R$ large, for Lemmas~\ref{lem:Ufan},
\ref{lem:pos}, \ref{lem:regions}, \ref{lem:chamber} and~\ref{lem:T0};
\item the leg collar paths and the slides, the arcs $\gamma_q$ and the extensions $\Theta$ of Lemma~\ref{lem:collar}
(Section~\ref{ssec:res-local});
\item the extension $F$ (Lemma~\ref{lem:ext});
\item the path in $\mathcal R$ (Lemma~\ref{lem:T}).
\end{enumerate}

The local pieces and Lemma~\ref{lem:ext} define a continuous map $F\colon\mathfrak D\to W$. Within each local domain the pieces are
continuous: the leg collar paths depend continuously on the point of the leg between switching points
(Section~\ref{ssec:res-local}), each switching point lies in a switching disc, and by Lemma~\ref{lem:collar} the maps $\Theta$
extend the collar paths; the circle actions used ($S^1_v$, and $T^2$ at the negative vertices) act exactly on the relevant parts of
$W$ and fix $\mathrm{Crit}$ where the domain collapses the circle. Different local domains meet along transverse segments at
points of the legs, where both are given by the same (slid) leg collar paths, and the local domains meet the regular part along
the inner boundary (I), which consists of $S^1_v$-orbits through points of $\mathfrak s$ and of the sets $\mathcal Q\cdot\mathfrak s(b)$,
where Lemma~\ref{lem:ext} extends them. By Lemma~\ref{lem:domain},
\[
\partial F_*[\mathfrak D]=\sum_qF_*[\partial_q\mathfrak D],
\]
and by Lemma~\ref{lem:nodedeg}, $F_*[\partial_q\mathfrak D]$ is a $2$-cycle supported on $C^W_q$ of degree $-\rho_q(S)$. A
$2$-cycle on $C^W_q\cong S^2$ of degree $d$ differs from $d$ times a fundamental cycle by the boundary of a $3$-chain in $C^W_q$;
adding these chains, we obtain a $3$-chain $\Xi_W$ on $W$, supported in $F(\mathfrak D)\cup\bigcup_qC^W_q$, with
\[
\Xi_W\equiv-F_*[\mathfrak D]\pmod{\text{chains in }\textstyle\bigcup_qC^W_q},\qquad\partial\Xi_W=\sum_q\rho_q(S)[C^W_q].
\]
Put $\Xi(S)=(\Upsilon_{\rm res})_*\Xi_W$. By Lemmas~\ref{lem:T0} and~\ref{lem:T},
\[
\partial\,\Xi(S)=\sum_q\rho_q(S)\,[C_q].
\]
The chain is constructed from $(S,j,v)$ and the listed choices, and no relation among the $[C_q]$ is used. This proves the first
assertion of Theorem~\ref{thm:res}. In particular $\sum_q\rho_q(S)[C_q]=0$ in $H_2(\widehat X;\Z)$, hence in
$H_2(\widehat X;\Q)$, and $\rho(S)\in K$ by Proposition~\ref{prop:partialres}.

The node coefficients in the first assertion are those of the diagonals of $\widehat X$. With the diagonals fixed in
Section~\ref{ssec:admissible}, which need not be those of $\widehat X$, the identity holds after reversing the sign of $\rho_q$
at every node where the two differ, since a change of diagonal reverses the sign of $\rho_q$ and not that of $[C_q]$. The
second assertion of Theorem~\ref{thm:res}, for an arbitrary small
resolution, follows from Lemma~\ref{lem:flop}, applied to $\lambda=\rho(S)$ and the chain $\Xi(S)$.

\begin{lemma}[flops]\label{lem:flop}
Let $\widehat X$ and $\widehat X'$ be small resolutions of $X'_\sigma$ whose diagonals differ exactly at the nodes of a set
$Q'\subseteq Q$, with exceptional curves $C_q$ and $C'_q$, and let $\lambda\in\Z^Q$. Put $\lambda'_q=-\lambda_q$ for $q\in Q'$ and
$\lambda'_q=\lambda_q$ otherwise. Then $\sum_q\lambda_q[C_q]=0$ in $H_2(\widehat X;\Z)$ if and only if
$\sum_q\lambda'_q[C'_q]=0$ in $H_2(\widehat X';\Z)$. More precisely, the image of a $3$-chain $\Xi$ on $\widehat X$ with
$\partial\Xi=\sum_q\lambda_q[C_q]$ in $H_3(\widehat X,\bigcup_q\widehat V_q)\cong H_3(X^\circ,\bigcup_q\partial V_q)$ is also the
class of a $3$-chain $\Xi'$ on $\widehat X'$ with $\partial\Xi'=\sum_q\lambda'_q[C'_q]$, where $\widehat V_q$, $\partial V_q$ and
$X^\circ$ are as in the proof.
\end{lemma}
\begin{proof}
Let $V_q\subset X'_\sigma$ be small disjoint closed conical neighbourhoods of the nodes, with boundaries
$\partial V_q\cong S^2\times S^3$, and let $\widehat V_q\subset\widehat X$ and $\widehat V'_q\subset\widehat X'$ be their
preimages; these are tubular neighbourhoods of $C_q$ and $C'_q$, with boundary $\partial V_q$. Both resolutions are isomorphisms
over $X^\circ:=X'_\sigma\setminus\bigcup_q\operatorname{int}V_q$. By excision,
$H_3(\widehat X,\bigcup_q\widehat V_q)\cong H_3(X^\circ,\bigcup_q\partial V_q)\cong H_3(\widehat X',\bigcup_q\widehat V'_q)$, and
the connecting maps of the two pairs factor as
\[
H_3(X^\circ,\textstyle\bigcup_q\partial V_q)\xrightarrow{\ \partial\ }\bigoplus_qH_2(\partial V_q)\longrightarrow
\bigoplus_qH_2(\widehat V_q),\ \text{respectively}\ \bigoplus_qH_2(\widehat V'_q),
\]
where $H_2(\partial V_q)\cong\Z$ maps isomorphically onto $H_2(\widehat V_q)=\Z[C_q]$ and onto $H_2(\widehat V'_q)=\Z[C'_q]$.
Let $\ell_q\in H_2(\partial V_q)$ map to $[C_q]$. The class $D_b$ of Section~\ref{ssec:partialres}, $b$ a vertex of the node
parallelogram off the diagonal $\{a,c\}$ of $\widehat X$, is defined on $X^\circ$ and restricts to the same class on $\partial
V_q$ for both resolutions, and $\langle D_b,\ell_q\rangle=D_b\cdot C_q=1$. For $q\in Q'$ the curve $C'_q$ is the curve of the wall
$\operatorname{cone}(b,d)$, and the relation $a+c-b-d=0$ gives $D_b\cdot C'_q=-1$; so $\ell_q$ maps to $-[C'_q]$. For $q\notin Q'$,
$\widehat V'_q=\widehat V_q$ and $\ell_q$ maps to $[C'_q]=[C_q]$. By the exact sequences of the pairs,
$\sum_q\lambda_q[C_q]=0$ in $H_2(\widehat X)$ if and only if $\sum_q\lambda_q\ell_q$ lies in the image of $\partial$, if and only
if $\sum_q\lambda'_q[C'_q]=0$ in $H_2(\widehat X')$, the preimage under $\partial$ being the same class; this also gives the
statement about chains.
\end{proof}

A small resolution that is not projective is handled in the same way: by Proposition~\ref{prop:partialres}(3) and
Lemma~\ref{lem:regref} there is a projective small resolution induced by a regular refinement of $\Sigma'$, the chain is constructed there, and Lemma~\ref{lem:flop} carries it to
the given one. This persistence under a change of resolution agrees with \cite[Rem.~9.4]{CBM} for relations induced by
tropical $2$-cycles. Lemma~\ref{lem:flop} records the sign change at each flopped curve when every exceptional curve is
given its complex orientation.

\emph{The last assertion of Theorem~\ref{thm:res}.} Label the node parallelograms by the diagonals of $\widehat X$, and orient
the vanishing spheres $\delta_q$ by Definition~\ref{def:orient} with this labelling. Triangulated, $S$ is a relative tropical
$2$-cycle $z_S$ with node coefficients $\rho(S)$ (Lemma~\ref{lem:lift}), so for the members $Y_t$ of Theorem~\ref{thm:main}
Theorem~\ref{thm:lift} gives a $4$-chain $\check\Xi(S)$ with $\partial\check\Xi(S)=\sum_q\rho_q(S)[\delta_q]$. Together with the
first assertion,
\[
\sum_q\rho_q(S)\,[C_q]=0\ \text{ in }H_2(\widehat X;\Z),\qquad\sum_q\rho_q(S)\,[\delta_q]=0\ \text{ in }H_3(Y_t;\Z),
\]
with the same coefficients on both sides. This completes the proof of Theorem~\ref{thm:res}. \qed

\begin{remark}\label{rem:res-choices}
(1) The support of $\Xi(S)$ can be taken inside $\Upsilon_{\rm res}$ of the part of $W$ lying over the $\eta$-neighbourhood of
$\varpi^{-1}$ of an arbitrarily small neighbourhood of $\hat S$ in $\widehat B$, tentacles included: by Lemma~\ref{lem:AC}(3), and
because every local piece is a family of local fibres over points of $\hat S\cap\widehat\Gamma$.

(2) The relative class of $\Xi(S)$ is not claimed to be independent of the choices; the extensions $\Theta$ (for instance the
sign of a half-turn), the switching discs, the slides and the path in $\mathcal R$ may change it by closed $3$-cycles.

(3) The identity $\sum_q\rho_q(S)\circv_q=0$ is proved combinatorially in Proposition~\ref{prop:containment}, from a collar of
$\partial S$ alone. Theorem~\ref{thm:res} realises it by a chain built from all of $S$.
\end{remark}

\begin{remark}[root classes at the del Pezzo points]\label{rem:res-roots}
Let $G$ be a pentagon or a hexagon of $\Delta$, with the diagonals through its centre, and let $E$ be the exceptional surface
of $\widehat X$ over the singular point of $X$ on the orbit of $\operatorname{cone}(G)$ (Proposition~\ref{prop:classical}).
Since $\rho(S)\in K$, the part $\sum_{q\subset G}\rho_q(S)[C_q]$ of $\partial\Xi(S)$ in $E$ is a class in the root lattice of
Remark~\ref{rem:tie}: a multiple of $e_1-e_2$ at a pentagon, of $\ell$ at a hexagon on the $A_1$ branch, and an element of the
$A_2$ root lattice at a hexagon on the $A_2$ branch. So $\Xi(S)$ defines a relative class on the complement of small
neighbourhoods of the singular points of $X$, whose boundary at a del Pezzo point lies in that root lattice, and at an
ordinary node of $X$ is $\rho_q$ times the generator of $H_2(S^2\times S^3)$.
\end{remark}

\subsection{Relative tropical \texorpdfstring{$2$}{2}-cycles}\label{ssec:res-relative}
Theorem~\ref{thm:res} is stated for tropical $2$-cycles in the sense of \cite{CBM}. For the relative tropical $2$-cycles of
Definition~\ref{def:relcycle} the relation itself needs no chain: $\rho(z)\in K$ for every relative cycle $z$ on every
triangulation in which $\Gamma$ is full (Theorem~\ref{thm:real}), so $\sum_q\rho_q(z)[C_q]=0$ in $H_2(\widehat X;\Z)$ by
Proposition~\ref{prop:integral}.

\begin{remark}\label{rem:res-relative}
The construction of this section uses the strictness of $S$ mainly through the shape of the domain: one primitive field on each
sheet, trivalent interior edges and four-valent interior vertices, and a boundary that passes each node at most once, along one
pair of opposite legs. Its ingredients are nevertheless local: the pieces
along the legs, at the negative and positive vertices, at the crossings and at the nodes, and the extension over aspherical
regions. Lemma~\ref{lem:nodedeg} is already stated for arbitrary leg coefficients $c_{ab}$, $c_{bc}$, and the crossing pieces
$S^1_v\cdot\mathfrak s(D_k)$ are solid tori for every monodromy-invariant field. The mirror side removes the corresponding restrictions by building the regular part over an abstract domain and by
removing fibre-class boundary with the positive real locus (Section~\ref{sec:lift}, Proposition~\ref{lem:F}); on $\widehat X$ we
have no analogue of the second step, and we do not extend Theorem~\ref{thm:res} to relative tropical $2$-cycles here.
Nothing in this paper depends on such an extension.
\end{remark}
